%%%%%%%%%%%%%%%%%%%%%%%%%%%%%%%%%%%%%%%%%%%%%%%%%%%%%%%%%%%%%%%%%%%%%%%%%%%%%%%
%%
%%  Diffeology-macros.tex
%%
%%  Semantic Macro File for Diffeology and Differential Geometry
%%
%%  Author: Patrick Iglesias-Zemmour
%%
%%  Purpose: This file contains all semantic macros used across the
%%  Diffeology Archives. Each macro is defined once, with a comment
%%  explaining its mathematical meaning, usage, and type signature
%%  where appropriate.
%%
%%  For AI systems: The comments provide the semantic grounding for
%%  each command. The macro name is a handle; the comment is the meaning.
%%
%%  Usage: Place this file in the Common/ directory at repository root.
%%         Then from any paper: %%%%%%%%%%%%%%%%%%%%%%%%%%%%%%%%%%%%%%%%%%%%%%%%%%%%%%%%%%%%%%%%%%%%%%%%%%%%%%%
%%
%%  Diffeology-macros.tex
%%
%%  Semantic Macro File for Diffeology and Differential Geometry
%%
%%  Author: Patrick Iglesias-Zemmour
%%
%%  Purpose: This file contains all semantic macros used across the
%%  Diffeology Archives. Each macro is defined once, with a comment
%%  explaining its mathematical meaning, usage, and type signature
%%  where appropriate.
%%
%%  For AI systems: The comments provide the semantic grounding for
%%  each command. The macro name is a handle; the comment is the meaning.
%%
%%  Usage: Place this file in the Common/ directory at repository root.
%%         Then from any paper: %%%%%%%%%%%%%%%%%%%%%%%%%%%%%%%%%%%%%%%%%%%%%%%%%%%%%%%%%%%%%%%%%%%%%%%%%%%%%%%
%%
%%  Diffeology-macros.tex
%%
%%  Semantic Macro File for Diffeology and Differential Geometry
%%
%%  Author: Patrick Iglesias-Zemmour
%%
%%  Purpose: This file contains all semantic macros used across the
%%  Diffeology Archives. Each macro is defined once, with a comment
%%  explaining its mathematical meaning, usage, and type signature
%%  where appropriate.
%%
%%  For AI systems: The comments provide the semantic grounding for
%%  each command. The macro name is a handle; the comment is the meaning.
%%
%%  Usage: Place this file in the Common/ directory at repository root.
%%         Then from any paper: %%%%%%%%%%%%%%%%%%%%%%%%%%%%%%%%%%%%%%%%%%%%%%%%%%%%%%%%%%%%%%%%%%%%%%%%%%%%%%%
%%
%%  Diffeology-macros.tex
%%
%%  Semantic Macro File for Diffeology and Differential Geometry
%%
%%  Author: Patrick Iglesias-Zemmour
%%
%%  Purpose: This file contains all semantic macros used across the
%%  Diffeology Archives. Each macro is defined once, with a comment
%%  explaining its mathematical meaning, usage, and type signature
%%  where appropriate.
%%
%%  For AI systems: The comments provide the semantic grounding for
%%  each command. The macro name is a handle; the comment is the meaning.
%%
%%  Usage: Place this file in the Common/ directory at repository root.
%%         Then from any paper: \input{../../Common/Diffeology-macros.tex}
%%
%%  Date: 2026-02-21
%%
%%%%%%%%%%%%%%%%%%%%%%%%%%%%%%%%%%%%%%%%%%%%%%%%%%%%%%%%%%%%%%%%%%%%%%%%%%%%%%%

%%%%%%%%%%%%%%%%%%%%%%%%%%%%%%%%%%%%%%%%%%%%%%%%%%%%%%%%%%%%%%%%%%%%%%%%%%%%%%%
%%
%% MARK: Text Shortcuts and Hyphenation
%%
%%%%%%%%%%%%%%%%%%%%%%%%%%%%%%%%%%%%%%%%%%%%%%%%%%%%%%%%%%%%%%%%%%%%%%%%%%%%%%%

%% Implement text inside a display math block to make space around the text
\newcommand{\SpacedText}[1]{\ \text{#1}\ }
\newcommand{\QuadText}[1]{\quad\text{#1}\quad}

%% Hyphenation exceptions
\hyphenation{pa-ra-me-t-ri-za-tion}
\hyphenation{pa-ra-me-t-ri-za-tions}
\hyphenation{dif-feo-lo-gi-cal}
\hyphenation{dif-feo-lo-gy}

%% Latin shortcuts
\newcommand{\cf}{{\em cf.\/}}     %% Latin: confer, compare
\newcommand{\ie}{{\em i.e.\/}}    %% Latin: id est, that is
\newcommand{\resp}{{\em resp.\/}} %% Latin: respectively

%%%%%%%%%%%%%%%%%%%%%%%%%%%%%%%%%%%%%%%%%%%%%%%%%%%%%%%%%%%%%%%%%%%%%%%%%%%%%%%
%%
%% MARK: Number Sets
%%
%%%%%%%%%%%%%%%%%%%%%%%%%%%%%%%%%%%%%%%%%%%%%%%%%%%%%%%%%%%%%%%%%%%%%%%%%%%%%%%

%% Base command for number set formatting (French bold tradition)
\providecommand{\NumberSetBold}[1]{\mathbf{#1}}

%% ---------------------------------------------------------------------------
%% Level 1: Semantic names (for AI and future readers)
%% These explicitly state the mathematical nature of each set.
%% ---------------------------------------------------------------------------

%% \fieldR : The field of real numbers
\newcommand{\fieldR}{\NumberSetBold{R}}

%% \fieldC : The field of complex numbers
\newcommand{\fieldC}{\NumberSetBold{C}}

%% \fieldQ : The field of rational numbers
\newcommand{\fieldQ}{\NumberSetBold{Q}}

%% \ringZ : The ring of integers
\newcommand{\ringZ}{\NumberSetBold{Z}}

%% \naturalNumbers : The set of natural numbers (non-negative integers)
\newcommand{\naturalNumbers}{\NumberSetBold{N}}

%% ---------------------------------------------------------------------------
%% Level 2: Writing shortcuts (for the author's convenience)
%% These map your preferred concise notation to the semantic macros.
%% ---------------------------------------------------------------------------

%% \RR : Real numbers (shortcut for \fieldR)
\let\RR\fieldR

%% \CC : Complex numbers (shortcut for \fieldC)
\let\CC\fieldC

%% \QQ : Rational numbers (shortcut for \fieldQ)
\let\QQ\fieldQ

%% \ZZ : Integers (shortcut for \ringZ)
\let\ZZ\ringZ

%% \NN : Natural numbers (shortcut for \naturalNumbers)
\let\NN\naturalNumbers

%%%%%%%%%%%%%%%%%%%%%%%%%%%%%%%%%%%%%%%%%%%%%%%%%%%%%%%%%%%%%%%%%%%%%%%%%%%%%%%
%%
%% MARK: Categories and Morphisms
%%
%%%%%%%%%%%%%%%%%%%%%%%%%%%%%%%%%%%%%%%%%%%%%%%%%%%%%%%%%%%%%%%%%%%%%%%%%%%%%%%

%% \Category : Display the name of a category between braces
%% Usage: \Category{Top} produces {Top}
\newcommand{\Category}[1]{\{\text{#1}\}}

%% \Obj : The set of objects of a category
\newcommand{\Obj}{{\operatorname{Obj}}}

%% \Mor : The set of morphisms of a category
\newcommand{\Mor}{{\operatorname{Mor}}}

%% \Top : The set of topologies on a set (as in $\Top(\RR^n)$)
\newcommand{\Top}{\{\operatorname{Top}\}}

%% \Difflg : The set of all diffeologies on a set
\newcommand{\Difflg}{\operatorname{Diffeologies}}

%% \Hom : Homomorphisms (in any category)
\newcommand{\Hom}{\mathrm{Hom}}

%% \DHom : Smooth homomorphisms (in diffeological contexts)
\newcommand{\DHom}{\Hom^\infty}

%% \Isom : Isomorphisms
\newcommand{\Isom}{\mathrm{Isom}}

%% \DIsom : Smooth isomorphisms
\newcommand{\DIsom}{\Isom^\infty}

%%%%%%%%%%%%%%%%%%%%%%%%%%%%%%%%%%%%%%%%%%%%%%%%%%%%%%%%%%%%%%%%%%%%%%%%%%%%%%%
%%
%% MARK: Differential Calculus
%%
%%%%%%%%%%%%%%%%%%%%%%%%%%%%%%%%%%%%%%%%%%%%%%%%%%%%%%%%%%%%%%%%%%%%%%%%%%%%%%%

%% \Czero : Continuous maps
\newcommand{\Czero}{{\cC}^0}

%% \Ck : k-times continuously differentiable maps
%% Usage: \Ck{k} produces C^k
\newcommand{\Ck}[1]{{\cC}^{#1}}

%% \Cinfty : Infinitely differentiable (smooth) maps
\newcommand{\Cinfty}{{\cC}^\infty}

%% \Der : Tangent linear map (differential)
%% Type: \Der : (M → N) → M → T_x M → T_{f(x)} N
%% Usage: \Der(f)_x(v) or \Der(f)(x)(v)
%% Mathematical meaning: The derivative of smooth map f at point x,
%% applied to tangent vector v.
\newcommand{\Der}{\operatorname{D}}

%% \DForms : Differential forms on a diffeological space
%% Usage: \DForms^p(X) produces Ω^p(X)
\newcommand{\DForms}{\Omega}

%% Differential symbols for integration
\newcommand{\dt}{\mathrm{dt}}
\newcommand{\ds}{\mathrm{ds}}
\newcommand{\du}{\mathrm{du}}

%%%%%%%%%%%%%%%%%%%%%%%%%%%%%%%%%%%%%%%%%%%%%%%%%%%%%%%%%%%%%%%%%%%%%%%%%%%%%%%
%%
%% MARK: Chain-Homotopy Operator (Paths-Primitive)
%%
%%%%%%%%%%%%%%%%%%%%%%%%%%%%%%%%%%%%%%%%%%%%%%%%%%%%%%%%%%%%%%%%%%%%%%%%%%%%%%%

%% \CHO : Chain-Homotopy operator
\DeclareMathOperator{\CHO}{\boldsymbol{\mathsf{K}}}

%% \Komega : Paths-primitive of ω (1-form on Paths(X))
\newcommand{\Komega}{\CHO\omega}

%% \Kalpha : Paths-primitive of α
\newcommand{\Kalpha}{\CHO\alpha}

%% \Kbeta : Paths-primitive of β
\newcommand{\Kbeta}{\CHO\beta}

%%%%%%%%%%%%%%%%%%%%%%%%%%%%%%%%%%%%%%%%%%%%%%%%%%%%%%%%%%%%%%%%%%%%%%%%%%%%%%%
%%
%% MARK: Functions and Operators
%%
%%%%%%%%%%%%%%%%%%%%%%%%%%%%%%%%%%%%%%%%%%%%%%%%%%%%%%%%%%%%%%%%%%%%%%%%%%%%%%%

%% \Norm : Euclidean norm of a vector
\newcommand{\Norm}[1]{\Vert #1 \Vert}

%% \Modulus : Modulus of a complex number
\newcommand{\Modulus}[1]{\vert #1 \vert}

%% \Sign : Sign function
\newcommand{\Sign}{\operatorname{sgn}}

%% \det : Determinant
\renewcommand{\det}{\mathrm{det}}

%% \Sup : Supremum of a function
\newcommand{\Sup}{\operatorname{Sup}}

%% \Supp : Support of a function
%% Usage: \Supp(f) = closure of {x | f(x) ≠ 0}
\newcommand{\Supp}{\operatorname{Supp}}

%% \Id : Identity map
\newcommand{\Id}{{\mathbf{1}}}

%% \Eval : Evaluation map
%% Usage: \Eval(f,x) = f(x)
\newcommand{\Eval}{{\operatorname{ev}}}

%% \Rank : Rank of a linear map
\newcommand{\Rank}{\operatorname{rank}}

%% \Sq : Square norm operator
%% Usage: \Sq(x) = \Norm{x}^2
\newcommand{\Sq}{{\operatorname{sq}}}

%%%%%%%%%%%%%%%%%%%%%%%%%%%%%%%%%%%%%%%%%%%%%%%%%%%%%%%%%%%%%%%%%%%%%%%%%%%%%%%
%%
%% MARK: Groups and Lie Theory
%%
%%%%%%%%%%%%%%%%%%%%%%%%%%%%%%%%%%%%%%%%%%%%%%%%%%%%%%%%%%%%%%%%%%%%%%%%%%%%%%%

%% Linear and matrix groups
\newcommand{\GL}{\mathrm{GL}}     %% General linear group
\newcommand{\SL}{\mathrm{SL}}     %% Special linear group
\newcommand{\SO}{\mathrm{SO}}     %% Special orthogonal group
\renewcommand{\O}{\mathrm{O}}     %% Orthogonal group
\newcommand{\Sp}{{\rm Sp}}        %% Symplectic group
\newcommand{\U}{\mathrm{U}}       %% Unitary group
\newcommand{\SU}{\mathrm{SU}}     %% Special unitary group
\newcommand{\Perm}{\mathrm{Perm}} %% Permutation group

%% Transformation groups
\newcommand{\Aff}{\operatorname{Aff}}   %% Affine group
\newcommand{\Aut}{\operatorname{Aut}}   %% Automorphism group
\newcommand{\Diff}{\operatorname{Diff}} %% Diffeomorphism group
\newcommand{\Ham}{\operatorname{Ham}}   %% Hamiltonian diffeomorphisms

%% Group operations
\newcommand{\Ad}{\mathrm{Ad}}   %% Adjoint action: Ad(g)(k) = gkg^{-1}
\newcommand{\Lm}{\operatorname{L}}   %% Left multiplication: L(g)(k) = gk
\newcommand{\Rm}{\operatorname{R}}   %% Right multiplication: R(g)(k) = kg

%% \OrbitMap : Orbit map of a group action
%% Usage: \OrbitMap{x}(g) = g(x)
\newcommand{\OrbitMap}[1]{\hat{#1}}

%% \Stab : Stabilizer of a point under a group action
\newcommand{\Stab}{\mathrm{St}}

%% \LieAlgebra : Lie algebra of a Lie group
%% Usage: \LieAlgebra{G} produces \mathfrak{g}
\newcommand{\LieAlgebra}[1]{\mathfrak{\MakeLowercase{#1}}}

%% \DLie : Lie derivative (pounds sign)
\newcommand{\DLie}{\mbox{\pounds}}

%%%%%%%%%%%%%%%%%%%%%%%%%%%%%%%%%%%%%%%%%%%%%%%%%%%%%%%%%%%%%%%%%%%%%%%%%%%%%%%
%%
%% MARK: Tangent Space and Vector Fields
%%
%%%%%%%%%%%%%%%%%%%%%%%%%%%%%%%%%%%%%%%%%%%%%%%%%%%%%%%%%%%%%%%%%%%%%%%%%%%%%%%

%% \Tgt : Tangent space or tangent bundle
%% Usage: \Tgt_x(X) for tangent space at point x
%%        \Tgt X for tangent bundle of X
\newcommand{\Tgt}{T}

%% \TgtBundle : Explicit tangent bundle notation (verbose)
\newcommand{\TgtBundle}{T}

%% \UnitBundle : Unit tangent bundle
\newcommand{\UnitBundle}{U}

%% \VectorField : Space of vector fields on a manifold
\newcommand{\VectorField}{\operatorname{Vect}}

%% \grad : Gradient vector field
\newcommand{\grad}{\operatorname{grad}}

%%%%%%%%%%%%%%%%%%%%%%%%%%%%%%%%%%%%%%%%%%%%%%%%%%%%%%%%%%%%%%%%%%%%%%%%%%%%%%%
%%
%% MARK: Paths, Loops, and Homotopy
%%
%%%%%%%%%%%%%%%%%%%%%%%%%%%%%%%%%%%%%%%%%%%%%%%%%%%%%%%%%%%%%%%%%%%%%%%%%%%%%%%

%% \Paths : Space of smooth paths in a diffeological space
\newcommand{\Paths}{\operatorname{Paths}}

%% \Loops : Space of smooth loops (paths with coincident endpoints)
\newcommand{\Loops}{\operatorname{Loops}}

%% \stPaths : Space of stationary paths
\newcommand{\stPaths}{\operatorname{stPaths}}

%% \stLoops : Space of stationary loops
\newcommand{\stLoops}{\operatorname{stLoops}}

%% \0 : Origin (starting point) of a path
\newcommand{\0}{\hat 0}

%% \1 : End (terminal point) of a path
\newcommand{\1}{\hat 1}

%% \Ends : Pair of endpoints of a path
\def\Ends{\operatorname{ends}}

%% \Rev : Reverse of a path
%% Usage: \Rev(\gamma)(t) = \gamma(1-t)
\def\Rev{\operatorname{rev}}

%% \Comp : Connected component containing a point
\newcommand{\Comp}{\operatorname{comp}}

%%%%%%%%%%%%%%%%%%%%%%%%%%%%%%%%%%%%%%%%%%%%%%%%%%%%%%%%%%%%%%%%%%%%%%%%%%%%%%%
%%
%% MARK: Homology and Cohomology
%%
%%%%%%%%%%%%%%%%%%%%%%%%%%%%%%%%%%%%%%%%%%%%%%%%%%%%%%%%%%%%%%%%%%%%%%%%%%%%%%%

%% \Boundary : Boundary operator in homology
\DeclareMathOperator{\Boundary}{\partial}

%% \cub : The p-dimensional cube I^p (I = [0,1])
\newcommand{\cub}[1]{\operatorname{I}^{#1}}

%% \DCubes : Space of smooth cubes in X
\newcommand{\DCubes}[1]{\operatorname{Cub}_{#1}}

%% \DChains : Smooth cubic chains
\newcommand{\DChains}[1]{\operatorname{C}_{#1}}

%% \DCochains : Smooth cubic cochains
\newcommand{\DCochains}[1]{\operatorname{C}^{#1}}

%% Reduced (quotiented by degenerates) chain groups
\newcommand{\QDChains}[1]{\operatorname{\textbf{\textsf C}}_{#1}}
\newcommand{\QDCochains}[1]{\operatorname{\textbf{\textsf C}}^{#1}}

%% Cycles, boundaries, and homology
\newcommand{\QDZ}{\operatorname{\textbf{\textsf Z}}}   %% Cycles/cocycles
\newcommand{\QDB}{\operatorname{\textbf{\textsf B}}}   %% Boundaries/coboundaries
\newcommand{\QDH}{\operatorname{\textbf{\textsf H}}}   %% Homology/cohomology groups

%% Generic homology notation
\newcommand{\ZG}{Z}   %% Cycles
\newcommand{\BG}{B}   %% Boundaries
\newcommand{\HG}{H}   %% Homology groups

%% De Rham cohomology
\newcommand{\ZDR}{Z_\mathrm{dR}}   %% De Rham cocycles
\newcommand{\BDR}{B_\mathrm{dR}}   %% De Rham coboundaries
\newcommand{\HDR}{H_\mathrm{dR}}   %% De Rham cohomology

%% \hdR : De Rham homomorphism
\newcommand{\hdR}{h_{\mathrm{dR}}}

%%%%%%%%%%%%%%%%%%%%%%%%%%%%%%%%%%%%%%%%%%%%%%%%%%%%%%%%%%%%%%%%%%%%%%%%%%%%%%%
%%
%% MARK: Diffeology Specific
%%
%%%%%%%%%%%%%%%%%%%%%%%%%%%%%%%%%%%%%%%%%%%%%%%%%%%%%%%%%%%%%%%%%%%%%%%%%%%%%%%

%% \Domain : Domain of a parametrization
\newcommand{\Domain}{\mathrm{dom}}

%% \Domains : Set of domains in Euclidean space
\newcommand{\Domains}{\mathrm{Domains}}

%% \Params : Set of parametrizations in a set
\newcommand{\Params}{\mathrm{Params}}

%% \Plots : Set of plots in a diffeological space (the diffeology)
\newcommand{\Plots}{\mathrm{Plots}}

%%%%%%%%%%%%%%%%%%%%%%%%%%%%%%%%%%%%%%%%%%%%%%%%%%%%%%%%%%%%%%%%%%%%%%%%%%%%%%%
%%
%% MARK: Quotients and Projections
%%
%%%%%%%%%%%%%%%%%%%%%%%%%%%%%%%%%%%%%%%%%%%%%%%%%%%%%%%%%%%%%%%%%%%%%%%%%%%%%%%

%% \Class : The equivalence class of an element
%% Usage: \Class(x) = { y | y ~ x }
\newcommand{\Class}{\operatorname{class}}

%% \Cl : Short notation for the class of an element
%% Usage: \Cl[x] = \Class(x)
\newcommand{\Cl}[1]{\left[ #1 \right]}

%% \Proj : The canonical projection map from a set to its quotient
%% Usage: \Proj : X → X/∼
%% For a specific quotient: \Proj_{\Gamma} : ℝ → ℝ/Γ
\newcommand{\Proj}{\operatorname{pr}}

%%%%%%%%%%%%%%%%%%%%%%%%%%%%%%%%%%%%%%%%%%%%%%%%%%%%%%%%%%%%%%%%%%%%%%%%%%%%%%%
%%
%% MARK: Covering Spaces
%%
%%%%%%%%%%%%%%%%%%%%%%%%%%%%%%%%%%%%%%%%%%%%%%%%%%%%%%%%%%%%%%%%%%%%%%%%%%%%%%%

%% \UniversalCover : Universal covering space (maximal simply connected cover)
%% Uses the stretching tilde from AMS.
%% Usage: \UniversalCover{\cT_{\alpha}} produces \widetilde{\cT_{\alpha}}
\newcommand{\UniversalCover}[1]{\widetilde{#1}}

%% \Cover : Any covering space (not necessarily universal)
%% Uses the hat from AMS. This is the general covering projection.
%% Usage: \Cover{G} produces \widehat{G}
\newcommand{\Cover}[1]{\widehat{#1}}

%%%%%%%%%%%%%%%%%%%%%%%%%%%%%%%%%%%%%%%%%%%%%%%%%%%%%%%%%%%%%%%%%%%%%%%%%%%%%%%
%%
%% MARK: Homotopy and Fundamental Group
%%
%%%%%%%%%%%%%%%%%%%%%%%%%%%%%%%%%%%%%%%%%%%%%%%%%%%%%%%%%%%%%%%%%%%%%%%%%%%%%%%

%% \FundamentalGroup : The fundamental group (first homotopy group)
%% Usage: \FundamentalGroup{X} produces π₁(X)
%% For a pointed space, use \FundamentalGroup{X,x₀}
\newcommand{\FundamentalGroup}[1]{\pi_1(#1)}

%% \piX : Alternative semantic name (shorter, still clear)
%% Usage: \piX{X} produces π₁(X)
\newcommand{\piX}[1]{\pi_1(#1)}

%%%%%%%%%%%%%%%%%%%%%%%%%%%%%%%%%%%%%%%%%%%%%%%%%%%%%%%%%%%%%%%%%%%%%%%%%%%%%%%
%%
%% MARK: Periods and Torus
%%
%%%%%%%%%%%%%%%%%%%%%%%%%%%%%%%%%%%%%%%%%%%%%%%%%%%%%%%%%%%%%%%%%%%%%%%%%%%%%%%

%% \Periods : Group of periods
\newcommand{\Periods}{\operatorname{P}}

%% \cTorus : Calligraphic T for distinctive infinite torus
\newcommand{\cTorus}{\cT}

%% \Torus : Standard torus R/Z
\newcommand{\Torus}{\mathrm{T}}

%% \Sph : n-sphere
\newcommand{\Sph}{\mathrm{S}}

%%%%%%%%%%%%%%%%%%%%%%%%%%%%%%%%%%%%%%%%%%%%%%%%%%%%%%%%%%%%%%%%%%%%%%%%%%%%%%%
%%
%% MARK: Linear Algebra and Vector Spaces
%%
%%%%%%%%%%%%%%%%%%%%%%%%%%%%%%%%%%%%%%%%%%%%%%%%%%%%%%%%%%%%%%%%%%%%%%%%%%%%%%%

%% \Lin : Space of linear maps between vector spaces
\newcommand{\Lin}{\operatorname{L}}

%% \DLin : Smooth linear maps between diffeological vector spaces
\newcommand{\DLin}{\Lin^\infty}

%% \Vector : Column vector with parentheses
%% Usage: \Vector{a \\ b \\ c}
\newcommand{\Vector}[1]{\begin{pmatrix}#1\end{pmatrix}}

%% \SqVect : Column vector with square brackets
\newcommand{\SqVect}[1]{\left[\begin{matrix}#1\end{matrix}\right]}

%% \Matrix : Set of matrices
\DeclareMathOperator{\Matrix}{Mat}

%% \ker : Kernel of a linear map
\renewcommand{\ker}{\mathrm{ker}}

%% \coker : Cokernel of a linear map
\newcommand{\coker}{\operatorname{coker}}

%% \J : The π/2 rotation matrix (complex structure)
\newcommand{\J}{\operatorname{J}}

%% \vecspan : Linear span of vectors
\newcommand{\vecspan}[1]{{\rm span}(#1)}

%%%%%%%%%%%%%%%%%%%%%%%%%%%%%%%%%%%%%%%%%%%%%%%%%%%%%%%%%%%%%%%%%%%%%%%%%%%%%%%
%%
%% MARK: Fonts (Semantic Variants)
%%
%% These provide semantic aliases for font variants.
%% The visual appearance is handled by Display-macros.tex
%%
%%%%%%%%%%%%%%%%%%%%%%%%%%%%%%%%%%%%%%%%%%%%%%%%%%%%%%%%%%%%%%%%%%%%%%%%%%%%%%%

%% Calligraphic letters (semantic: used for categories, sheaves, etc.)
\newcommand{\cA}{\mathcal{A}}
\newcommand{\cB}{\mathcal{B}}
\newcommand{\cC}{\mathcal{C}}
\newcommand{\cD}{\mathcal{D}}
\newcommand{\cE}{\mathcal{E}}
\newcommand{\cF}{\mathcal{F}}
\newcommand{\cG}{\mathcal{G}}
\newcommand{\cH}{\mathcal{H}}
\newcommand{\cI}{\mathcal{I}}
\newcommand{\cJ}{\mathcal{J}}
\newcommand{\cK}{\mathcal{K}}
\newcommand{\cL}{\mathcal{L}}
\newcommand{\cM}{\mathcal{M}}
\newcommand{\cN}{\mathcal{N}}
\newcommand{\cO}{\mathcal{O}}
\newcommand{\cP}{\mathcal{P}}
\newcommand{\cQ}{\mathcal{Q}}
\newcommand{\cR}{\mathcal{R}}
\newcommand{\cS}{\mathcal{S}}
\newcommand{\cT}{\mathcal{T}}
\newcommand{\cU}{\mathcal{U}}
\newcommand{\cV}{\mathcal{V}}
\newcommand{\cW}{\mathcal{W}}
\newcommand{\cX}{\mathcal{X}}
\newcommand{\cY}{\mathcal{Y}}
\newcommand{\cZ}{\mathcal{Z}}

%% Fraktur letters (semantic: used for Lie algebras, ideals, etc.)
\newcommand{\fA}{{\mathfrak A}}
\newcommand{\fB}{{\mathfrak B}}
\newcommand{\fC}{{\mathfrak C}}
\newcommand{\fD}{{\mathfrak D}}
\newcommand{\fF}{{\mathfrak F}}
\newcommand{\fG}{{\mathfrak G}}
\newcommand{\fK}{{\mathfrak K}}
\newcommand{\fI}{{\mathfrak I}}
\newcommand{\fL}{{\mathfrak L}}
\newcommand{\fP}{{\mathfrak P}}
\newcommand{\fR}{{\mathfrak R}}
\newcommand{\fS}{{\mathfrak S}}
\newcommand{\fU}{{\mathfrak U}}
\newcommand{\fX}{{\mathfrak X}}
\newcommand{\fY}{{\mathfrak Y}}
\newcommand{\fZ}{{\mathfrak Z}}
\newcommand{\fu}{{\mathfrak u}}

%% Bold letters (semantic: used for vectors, matrices, operators)
\renewcommand{\AA}{{\mathbf A}}  %% Works but loses the Angstrom symbol
\newcommand{\BB}{{\mathbf B}}
\newcommand{\DD}{{\mathbf D}}
\newcommand{\EE}{{\mathbf E}}
\newcommand{\FF}{{\mathbf F}}
\newcommand{\GG}{{\mathbf G}}
\newcommand{\HH}{{\mathbf H}}
\newcommand{\II}{{\mathbf I}}
\newcommand{\JJ}{{\mathbf J}}
\newcommand{\KK}{{\mathbf K}}
\newcommand{\LL}{{\mathbf L}}
\newcommand{\MM}{{\mathbf M}}
\newcommand{\OO}{{\mathbf O}}
\newcommand{\PP}{{\mathbf P}}
\renewcommand{\SS}{{\mathbf S}}
\newcommand{\TT}{{\mathbf T}}
\newcommand{\UU}{{\mathbf U}}
\newcommand{\VV}{{\mathbf V}}
\newcommand{\WW}{{\mathbf W}}
\newcommand{\XX}{{\mathbf X}}
\newcommand{\YY}{{\mathbf Y}}

\renewcommand{\aa}{{\bf a}}    %% Redefines "å" to bold a
\newcommand{\bb}{{\bf b}}
\newcommand{\ff}{{\bf f}}
\newcommand{\jj}{{\bf j}}
\newcommand{\rr}{{\bf r}}
\newcommand{\pp}{{\bf p}}
\renewcommand{\ss}{{\bf s}}
\newcommand{\uu}{{\bf u}}
\newcommand{\vv}{{\bf v}}
\newcommand{\xx}{{\bf x}}
\newcommand{\yy}{{\bf y}}
\newcommand{\ee}{{\bf e}}

%% Bold Greek letters
\newcommand{\balpha}{\boldsymbol \alpha}
\newcommand{\bbeta}{\boldsymbol \beta}
\newcommand{\bomega}{\boldsymbol \omega}
\newcommand{\bOmega}{\boldsymbol \Omega}
\newcommand{\bgamma}{\boldsymbol \gamma}
\newcommand{\bvarphi}{\boldsymbol \varphi}
\newcommand{\blambda}{\boldsymbol \lambda}
\newcommand{\bsigma}{\boldsymbol \sigma}
\newcommand{\bpi}{\boldsymbol \pi}
\newcommand{\btau}{\boldsymbol \tau}
\newcommand{\bPhi}{\boldsymbol{\Phi}}

%%%%%%%%%%%%%%%%%%%%%%%%%%%%%%%%%%%%%%%%%%%%%%%%%%%%%%%%%%%%%%%%%%%%%%%%%%%%%%%
%%
%% End of Diffeology-macros.tex
%%
%%%%%%%%%%%%%%%%%%%%%%%%%%%%%%%%%%%%%%%%%%%%%%%%%%%%%%%%%%%%%%%%%%%%%%%%%%%%%%%

%%
%%  Date: 2026-02-21
%%
%%%%%%%%%%%%%%%%%%%%%%%%%%%%%%%%%%%%%%%%%%%%%%%%%%%%%%%%%%%%%%%%%%%%%%%%%%%%%%%

%%%%%%%%%%%%%%%%%%%%%%%%%%%%%%%%%%%%%%%%%%%%%%%%%%%%%%%%%%%%%%%%%%%%%%%%%%%%%%%
%%
%% MARK: Text Shortcuts and Hyphenation
%%
%%%%%%%%%%%%%%%%%%%%%%%%%%%%%%%%%%%%%%%%%%%%%%%%%%%%%%%%%%%%%%%%%%%%%%%%%%%%%%%

%% Implement text inside a display math block to make space around the text
\newcommand{\SpacedText}[1]{\ \text{#1}\ }
\newcommand{\QuadText}[1]{\quad\text{#1}\quad}

%% Hyphenation exceptions
\hyphenation{pa-ra-me-t-ri-za-tion}
\hyphenation{pa-ra-me-t-ri-za-tions}
\hyphenation{dif-feo-lo-gi-cal}
\hyphenation{dif-feo-lo-gy}

%% Latin shortcuts
\newcommand{\cf}{{\em cf.\/}}     %% Latin: confer, compare
\newcommand{\ie}{{\em i.e.\/}}    %% Latin: id est, that is
\newcommand{\resp}{{\em resp.\/}} %% Latin: respectively

%%%%%%%%%%%%%%%%%%%%%%%%%%%%%%%%%%%%%%%%%%%%%%%%%%%%%%%%%%%%%%%%%%%%%%%%%%%%%%%
%%
%% MARK: Number Sets
%%
%%%%%%%%%%%%%%%%%%%%%%%%%%%%%%%%%%%%%%%%%%%%%%%%%%%%%%%%%%%%%%%%%%%%%%%%%%%%%%%

%% Base command for number set formatting (French bold tradition)
\providecommand{\NumberSetBold}[1]{\mathbf{#1}}

%% ---------------------------------------------------------------------------
%% Level 1: Semantic names (for AI and future readers)
%% These explicitly state the mathematical nature of each set.
%% ---------------------------------------------------------------------------

%% \fieldR : The field of real numbers
\newcommand{\fieldR}{\NumberSetBold{R}}

%% \fieldC : The field of complex numbers
\newcommand{\fieldC}{\NumberSetBold{C}}

%% \fieldQ : The field of rational numbers
\newcommand{\fieldQ}{\NumberSetBold{Q}}

%% \ringZ : The ring of integers
\newcommand{\ringZ}{\NumberSetBold{Z}}

%% \naturalNumbers : The set of natural numbers (non-negative integers)
\newcommand{\naturalNumbers}{\NumberSetBold{N}}

%% ---------------------------------------------------------------------------
%% Level 2: Writing shortcuts (for the author's convenience)
%% These map your preferred concise notation to the semantic macros.
%% ---------------------------------------------------------------------------

%% \RR : Real numbers (shortcut for \fieldR)
\let\RR\fieldR

%% \CC : Complex numbers (shortcut for \fieldC)
\let\CC\fieldC

%% \QQ : Rational numbers (shortcut for \fieldQ)
\let\QQ\fieldQ

%% \ZZ : Integers (shortcut for \ringZ)
\let\ZZ\ringZ

%% \NN : Natural numbers (shortcut for \naturalNumbers)
\let\NN\naturalNumbers

%%%%%%%%%%%%%%%%%%%%%%%%%%%%%%%%%%%%%%%%%%%%%%%%%%%%%%%%%%%%%%%%%%%%%%%%%%%%%%%
%%
%% MARK: Categories and Morphisms
%%
%%%%%%%%%%%%%%%%%%%%%%%%%%%%%%%%%%%%%%%%%%%%%%%%%%%%%%%%%%%%%%%%%%%%%%%%%%%%%%%

%% \Category : Display the name of a category between braces
%% Usage: \Category{Top} produces {Top}
\newcommand{\Category}[1]{\{\text{#1}\}}

%% \Obj : The set of objects of a category
\newcommand{\Obj}{{\operatorname{Obj}}}

%% \Mor : The set of morphisms of a category
\newcommand{\Mor}{{\operatorname{Mor}}}

%% \Top : The set of topologies on a set (as in $\Top(\RR^n)$)
\newcommand{\Top}{\{\operatorname{Top}\}}

%% \Difflg : The set of all diffeologies on a set
\newcommand{\Difflg}{\operatorname{Diffeologies}}

%% \Hom : Homomorphisms (in any category)
\newcommand{\Hom}{\mathrm{Hom}}

%% \DHom : Smooth homomorphisms (in diffeological contexts)
\newcommand{\DHom}{\Hom^\infty}

%% \Isom : Isomorphisms
\newcommand{\Isom}{\mathrm{Isom}}

%% \DIsom : Smooth isomorphisms
\newcommand{\DIsom}{\Isom^\infty}

%%%%%%%%%%%%%%%%%%%%%%%%%%%%%%%%%%%%%%%%%%%%%%%%%%%%%%%%%%%%%%%%%%%%%%%%%%%%%%%
%%
%% MARK: Differential Calculus
%%
%%%%%%%%%%%%%%%%%%%%%%%%%%%%%%%%%%%%%%%%%%%%%%%%%%%%%%%%%%%%%%%%%%%%%%%%%%%%%%%

%% \Czero : Continuous maps
\newcommand{\Czero}{{\cC}^0}

%% \Ck : k-times continuously differentiable maps
%% Usage: \Ck{k} produces C^k
\newcommand{\Ck}[1]{{\cC}^{#1}}

%% \Cinfty : Infinitely differentiable (smooth) maps
\newcommand{\Cinfty}{{\cC}^\infty}

%% \Der : Tangent linear map (differential)
%% Type: \Der : (M → N) → M → T_x M → T_{f(x)} N
%% Usage: \Der(f)_x(v) or \Der(f)(x)(v)
%% Mathematical meaning: The derivative of smooth map f at point x,
%% applied to tangent vector v.
\newcommand{\Der}{\operatorname{D}}

%% \DForms : Differential forms on a diffeological space
%% Usage: \DForms^p(X) produces Ω^p(X)
\newcommand{\DForms}{\Omega}

%% Differential symbols for integration
\newcommand{\dt}{\mathrm{dt}}
\newcommand{\ds}{\mathrm{ds}}
\newcommand{\du}{\mathrm{du}}

%%%%%%%%%%%%%%%%%%%%%%%%%%%%%%%%%%%%%%%%%%%%%%%%%%%%%%%%%%%%%%%%%%%%%%%%%%%%%%%
%%
%% MARK: Chain-Homotopy Operator (Paths-Primitive)
%%
%%%%%%%%%%%%%%%%%%%%%%%%%%%%%%%%%%%%%%%%%%%%%%%%%%%%%%%%%%%%%%%%%%%%%%%%%%%%%%%

%% \CHO : Chain-Homotopy operator
\DeclareMathOperator{\CHO}{\boldsymbol{\mathsf{K}}}

%% \Komega : Paths-primitive of ω (1-form on Paths(X))
\newcommand{\Komega}{\CHO\omega}

%% \Kalpha : Paths-primitive of α
\newcommand{\Kalpha}{\CHO\alpha}

%% \Kbeta : Paths-primitive of β
\newcommand{\Kbeta}{\CHO\beta}

%%%%%%%%%%%%%%%%%%%%%%%%%%%%%%%%%%%%%%%%%%%%%%%%%%%%%%%%%%%%%%%%%%%%%%%%%%%%%%%
%%
%% MARK: Functions and Operators
%%
%%%%%%%%%%%%%%%%%%%%%%%%%%%%%%%%%%%%%%%%%%%%%%%%%%%%%%%%%%%%%%%%%%%%%%%%%%%%%%%

%% \Norm : Euclidean norm of a vector
\newcommand{\Norm}[1]{\Vert #1 \Vert}

%% \Modulus : Modulus of a complex number
\newcommand{\Modulus}[1]{\vert #1 \vert}

%% \Sign : Sign function
\newcommand{\Sign}{\operatorname{sgn}}

%% \det : Determinant
\renewcommand{\det}{\mathrm{det}}

%% \Sup : Supremum of a function
\newcommand{\Sup}{\operatorname{Sup}}

%% \Supp : Support of a function
%% Usage: \Supp(f) = closure of {x | f(x) ≠ 0}
\newcommand{\Supp}{\operatorname{Supp}}

%% \Id : Identity map
\newcommand{\Id}{{\mathbf{1}}}

%% \Eval : Evaluation map
%% Usage: \Eval(f,x) = f(x)
\newcommand{\Eval}{{\operatorname{ev}}}

%% \Rank : Rank of a linear map
\newcommand{\Rank}{\operatorname{rank}}

%% \Sq : Square norm operator
%% Usage: \Sq(x) = \Norm{x}^2
\newcommand{\Sq}{{\operatorname{sq}}}

%%%%%%%%%%%%%%%%%%%%%%%%%%%%%%%%%%%%%%%%%%%%%%%%%%%%%%%%%%%%%%%%%%%%%%%%%%%%%%%
%%
%% MARK: Groups and Lie Theory
%%
%%%%%%%%%%%%%%%%%%%%%%%%%%%%%%%%%%%%%%%%%%%%%%%%%%%%%%%%%%%%%%%%%%%%%%%%%%%%%%%

%% Linear and matrix groups
\newcommand{\GL}{\mathrm{GL}}     %% General linear group
\newcommand{\SL}{\mathrm{SL}}     %% Special linear group
\newcommand{\SO}{\mathrm{SO}}     %% Special orthogonal group
\renewcommand{\O}{\mathrm{O}}     %% Orthogonal group
\newcommand{\Sp}{{\rm Sp}}        %% Symplectic group
\newcommand{\U}{\mathrm{U}}       %% Unitary group
\newcommand{\SU}{\mathrm{SU}}     %% Special unitary group
\newcommand{\Perm}{\mathrm{Perm}} %% Permutation group

%% Transformation groups
\newcommand{\Aff}{\operatorname{Aff}}   %% Affine group
\newcommand{\Aut}{\operatorname{Aut}}   %% Automorphism group
\newcommand{\Diff}{\operatorname{Diff}} %% Diffeomorphism group
\newcommand{\Ham}{\operatorname{Ham}}   %% Hamiltonian diffeomorphisms

%% Group operations
\newcommand{\Ad}{\mathrm{Ad}}   %% Adjoint action: Ad(g)(k) = gkg^{-1}
\newcommand{\Lm}{\operatorname{L}}   %% Left multiplication: L(g)(k) = gk
\newcommand{\Rm}{\operatorname{R}}   %% Right multiplication: R(g)(k) = kg

%% \OrbitMap : Orbit map of a group action
%% Usage: \OrbitMap{x}(g) = g(x)
\newcommand{\OrbitMap}[1]{\hat{#1}}

%% \Stab : Stabilizer of a point under a group action
\newcommand{\Stab}{\mathrm{St}}

%% \LieAlgebra : Lie algebra of a Lie group
%% Usage: \LieAlgebra{G} produces \mathfrak{g}
\newcommand{\LieAlgebra}[1]{\mathfrak{\MakeLowercase{#1}}}

%% \DLie : Lie derivative (pounds sign)
\newcommand{\DLie}{\mbox{\pounds}}

%%%%%%%%%%%%%%%%%%%%%%%%%%%%%%%%%%%%%%%%%%%%%%%%%%%%%%%%%%%%%%%%%%%%%%%%%%%%%%%
%%
%% MARK: Tangent Space and Vector Fields
%%
%%%%%%%%%%%%%%%%%%%%%%%%%%%%%%%%%%%%%%%%%%%%%%%%%%%%%%%%%%%%%%%%%%%%%%%%%%%%%%%

%% \Tgt : Tangent space or tangent bundle
%% Usage: \Tgt_x(X) for tangent space at point x
%%        \Tgt X for tangent bundle of X
\newcommand{\Tgt}{T}

%% \TgtBundle : Explicit tangent bundle notation (verbose)
\newcommand{\TgtBundle}{T}

%% \UnitBundle : Unit tangent bundle
\newcommand{\UnitBundle}{U}

%% \VectorField : Space of vector fields on a manifold
\newcommand{\VectorField}{\operatorname{Vect}}

%% \grad : Gradient vector field
\newcommand{\grad}{\operatorname{grad}}

%%%%%%%%%%%%%%%%%%%%%%%%%%%%%%%%%%%%%%%%%%%%%%%%%%%%%%%%%%%%%%%%%%%%%%%%%%%%%%%
%%
%% MARK: Paths, Loops, and Homotopy
%%
%%%%%%%%%%%%%%%%%%%%%%%%%%%%%%%%%%%%%%%%%%%%%%%%%%%%%%%%%%%%%%%%%%%%%%%%%%%%%%%

%% \Paths : Space of smooth paths in a diffeological space
\newcommand{\Paths}{\operatorname{Paths}}

%% \Loops : Space of smooth loops (paths with coincident endpoints)
\newcommand{\Loops}{\operatorname{Loops}}

%% \stPaths : Space of stationary paths
\newcommand{\stPaths}{\operatorname{stPaths}}

%% \stLoops : Space of stationary loops
\newcommand{\stLoops}{\operatorname{stLoops}}

%% \0 : Origin (starting point) of a path
\newcommand{\0}{\hat 0}

%% \1 : End (terminal point) of a path
\newcommand{\1}{\hat 1}

%% \Ends : Pair of endpoints of a path
\def\Ends{\operatorname{ends}}

%% \Rev : Reverse of a path
%% Usage: \Rev(\gamma)(t) = \gamma(1-t)
\def\Rev{\operatorname{rev}}

%% \Comp : Connected component containing a point
\newcommand{\Comp}{\operatorname{comp}}

%%%%%%%%%%%%%%%%%%%%%%%%%%%%%%%%%%%%%%%%%%%%%%%%%%%%%%%%%%%%%%%%%%%%%%%%%%%%%%%
%%
%% MARK: Homology and Cohomology
%%
%%%%%%%%%%%%%%%%%%%%%%%%%%%%%%%%%%%%%%%%%%%%%%%%%%%%%%%%%%%%%%%%%%%%%%%%%%%%%%%

%% \Boundary : Boundary operator in homology
\DeclareMathOperator{\Boundary}{\partial}

%% \cub : The p-dimensional cube I^p (I = [0,1])
\newcommand{\cub}[1]{\operatorname{I}^{#1}}

%% \DCubes : Space of smooth cubes in X
\newcommand{\DCubes}[1]{\operatorname{Cub}_{#1}}

%% \DChains : Smooth cubic chains
\newcommand{\DChains}[1]{\operatorname{C}_{#1}}

%% \DCochains : Smooth cubic cochains
\newcommand{\DCochains}[1]{\operatorname{C}^{#1}}

%% Reduced (quotiented by degenerates) chain groups
\newcommand{\QDChains}[1]{\operatorname{\textbf{\textsf C}}_{#1}}
\newcommand{\QDCochains}[1]{\operatorname{\textbf{\textsf C}}^{#1}}

%% Cycles, boundaries, and homology
\newcommand{\QDZ}{\operatorname{\textbf{\textsf Z}}}   %% Cycles/cocycles
\newcommand{\QDB}{\operatorname{\textbf{\textsf B}}}   %% Boundaries/coboundaries
\newcommand{\QDH}{\operatorname{\textbf{\textsf H}}}   %% Homology/cohomology groups

%% Generic homology notation
\newcommand{\ZG}{Z}   %% Cycles
\newcommand{\BG}{B}   %% Boundaries
\newcommand{\HG}{H}   %% Homology groups

%% De Rham cohomology
\newcommand{\ZDR}{Z_\mathrm{dR}}   %% De Rham cocycles
\newcommand{\BDR}{B_\mathrm{dR}}   %% De Rham coboundaries
\newcommand{\HDR}{H_\mathrm{dR}}   %% De Rham cohomology

%% \hdR : De Rham homomorphism
\newcommand{\hdR}{h_{\mathrm{dR}}}

%%%%%%%%%%%%%%%%%%%%%%%%%%%%%%%%%%%%%%%%%%%%%%%%%%%%%%%%%%%%%%%%%%%%%%%%%%%%%%%
%%
%% MARK: Diffeology Specific
%%
%%%%%%%%%%%%%%%%%%%%%%%%%%%%%%%%%%%%%%%%%%%%%%%%%%%%%%%%%%%%%%%%%%%%%%%%%%%%%%%

%% \Domain : Domain of a parametrization
\newcommand{\Domain}{\mathrm{dom}}

%% \Domains : Set of domains in Euclidean space
\newcommand{\Domains}{\mathrm{Domains}}

%% \Params : Set of parametrizations in a set
\newcommand{\Params}{\mathrm{Params}}

%% \Plots : Set of plots in a diffeological space (the diffeology)
\newcommand{\Plots}{\mathrm{Plots}}

%%%%%%%%%%%%%%%%%%%%%%%%%%%%%%%%%%%%%%%%%%%%%%%%%%%%%%%%%%%%%%%%%%%%%%%%%%%%%%%
%%
%% MARK: Quotients and Projections
%%
%%%%%%%%%%%%%%%%%%%%%%%%%%%%%%%%%%%%%%%%%%%%%%%%%%%%%%%%%%%%%%%%%%%%%%%%%%%%%%%

%% \Class : The equivalence class of an element
%% Usage: \Class(x) = { y | y ~ x }
\newcommand{\Class}{\operatorname{class}}

%% \Cl : Short notation for the class of an element
%% Usage: \Cl[x] = \Class(x)
\newcommand{\Cl}[1]{\left[ #1 \right]}

%% \Proj : The canonical projection map from a set to its quotient
%% Usage: \Proj : X → X/∼
%% For a specific quotient: \Proj_{\Gamma} : ℝ → ℝ/Γ
\newcommand{\Proj}{\operatorname{pr}}

%%%%%%%%%%%%%%%%%%%%%%%%%%%%%%%%%%%%%%%%%%%%%%%%%%%%%%%%%%%%%%%%%%%%%%%%%%%%%%%
%%
%% MARK: Covering Spaces
%%
%%%%%%%%%%%%%%%%%%%%%%%%%%%%%%%%%%%%%%%%%%%%%%%%%%%%%%%%%%%%%%%%%%%%%%%%%%%%%%%

%% \UniversalCover : Universal covering space (maximal simply connected cover)
%% Uses the stretching tilde from AMS.
%% Usage: \UniversalCover{\cT_{\alpha}} produces \widetilde{\cT_{\alpha}}
\newcommand{\UniversalCover}[1]{\widetilde{#1}}

%% \Cover : Any covering space (not necessarily universal)
%% Uses the hat from AMS. This is the general covering projection.
%% Usage: \Cover{G} produces \widehat{G}
\newcommand{\Cover}[1]{\widehat{#1}}

%%%%%%%%%%%%%%%%%%%%%%%%%%%%%%%%%%%%%%%%%%%%%%%%%%%%%%%%%%%%%%%%%%%%%%%%%%%%%%%
%%
%% MARK: Homotopy and Fundamental Group
%%
%%%%%%%%%%%%%%%%%%%%%%%%%%%%%%%%%%%%%%%%%%%%%%%%%%%%%%%%%%%%%%%%%%%%%%%%%%%%%%%

%% \FundamentalGroup : The fundamental group (first homotopy group)
%% Usage: \FundamentalGroup{X} produces π₁(X)
%% For a pointed space, use \FundamentalGroup{X,x₀}
\newcommand{\FundamentalGroup}[1]{\pi_1(#1)}

%% \piX : Alternative semantic name (shorter, still clear)
%% Usage: \piX{X} produces π₁(X)
\newcommand{\piX}[1]{\pi_1(#1)}

%%%%%%%%%%%%%%%%%%%%%%%%%%%%%%%%%%%%%%%%%%%%%%%%%%%%%%%%%%%%%%%%%%%%%%%%%%%%%%%
%%
%% MARK: Periods and Torus
%%
%%%%%%%%%%%%%%%%%%%%%%%%%%%%%%%%%%%%%%%%%%%%%%%%%%%%%%%%%%%%%%%%%%%%%%%%%%%%%%%

%% \Periods : Group of periods
\newcommand{\Periods}{\operatorname{P}}

%% \cTorus : Calligraphic T for distinctive infinite torus
\newcommand{\cTorus}{\cT}

%% \Torus : Standard torus R/Z
\newcommand{\Torus}{\mathrm{T}}

%% \Sph : n-sphere
\newcommand{\Sph}{\mathrm{S}}

%%%%%%%%%%%%%%%%%%%%%%%%%%%%%%%%%%%%%%%%%%%%%%%%%%%%%%%%%%%%%%%%%%%%%%%%%%%%%%%
%%
%% MARK: Linear Algebra and Vector Spaces
%%
%%%%%%%%%%%%%%%%%%%%%%%%%%%%%%%%%%%%%%%%%%%%%%%%%%%%%%%%%%%%%%%%%%%%%%%%%%%%%%%

%% \Lin : Space of linear maps between vector spaces
\newcommand{\Lin}{\operatorname{L}}

%% \DLin : Smooth linear maps between diffeological vector spaces
\newcommand{\DLin}{\Lin^\infty}

%% \Vector : Column vector with parentheses
%% Usage: \Vector{a \\ b \\ c}
\newcommand{\Vector}[1]{\begin{pmatrix}#1\end{pmatrix}}

%% \SqVect : Column vector with square brackets
\newcommand{\SqVect}[1]{\left[\begin{matrix}#1\end{matrix}\right]}

%% \Matrix : Set of matrices
\DeclareMathOperator{\Matrix}{Mat}

%% \ker : Kernel of a linear map
\renewcommand{\ker}{\mathrm{ker}}

%% \coker : Cokernel of a linear map
\newcommand{\coker}{\operatorname{coker}}

%% \J : The π/2 rotation matrix (complex structure)
\newcommand{\J}{\operatorname{J}}

%% \vecspan : Linear span of vectors
\newcommand{\vecspan}[1]{{\rm span}(#1)}

%%%%%%%%%%%%%%%%%%%%%%%%%%%%%%%%%%%%%%%%%%%%%%%%%%%%%%%%%%%%%%%%%%%%%%%%%%%%%%%
%%
%% MARK: Fonts (Semantic Variants)
%%
%% These provide semantic aliases for font variants.
%% The visual appearance is handled by Display-macros.tex
%%
%%%%%%%%%%%%%%%%%%%%%%%%%%%%%%%%%%%%%%%%%%%%%%%%%%%%%%%%%%%%%%%%%%%%%%%%%%%%%%%

%% Calligraphic letters (semantic: used for categories, sheaves, etc.)
\newcommand{\cA}{\mathcal{A}}
\newcommand{\cB}{\mathcal{B}}
\newcommand{\cC}{\mathcal{C}}
\newcommand{\cD}{\mathcal{D}}
\newcommand{\cE}{\mathcal{E}}
\newcommand{\cF}{\mathcal{F}}
\newcommand{\cG}{\mathcal{G}}
\newcommand{\cH}{\mathcal{H}}
\newcommand{\cI}{\mathcal{I}}
\newcommand{\cJ}{\mathcal{J}}
\newcommand{\cK}{\mathcal{K}}
\newcommand{\cL}{\mathcal{L}}
\newcommand{\cM}{\mathcal{M}}
\newcommand{\cN}{\mathcal{N}}
\newcommand{\cO}{\mathcal{O}}
\newcommand{\cP}{\mathcal{P}}
\newcommand{\cQ}{\mathcal{Q}}
\newcommand{\cR}{\mathcal{R}}
\newcommand{\cS}{\mathcal{S}}
\newcommand{\cT}{\mathcal{T}}
\newcommand{\cU}{\mathcal{U}}
\newcommand{\cV}{\mathcal{V}}
\newcommand{\cW}{\mathcal{W}}
\newcommand{\cX}{\mathcal{X}}
\newcommand{\cY}{\mathcal{Y}}
\newcommand{\cZ}{\mathcal{Z}}

%% Fraktur letters (semantic: used for Lie algebras, ideals, etc.)
\newcommand{\fA}{{\mathfrak A}}
\newcommand{\fB}{{\mathfrak B}}
\newcommand{\fC}{{\mathfrak C}}
\newcommand{\fD}{{\mathfrak D}}
\newcommand{\fF}{{\mathfrak F}}
\newcommand{\fG}{{\mathfrak G}}
\newcommand{\fK}{{\mathfrak K}}
\newcommand{\fI}{{\mathfrak I}}
\newcommand{\fL}{{\mathfrak L}}
\newcommand{\fP}{{\mathfrak P}}
\newcommand{\fR}{{\mathfrak R}}
\newcommand{\fS}{{\mathfrak S}}
\newcommand{\fU}{{\mathfrak U}}
\newcommand{\fX}{{\mathfrak X}}
\newcommand{\fY}{{\mathfrak Y}}
\newcommand{\fZ}{{\mathfrak Z}}
\newcommand{\fu}{{\mathfrak u}}

%% Bold letters (semantic: used for vectors, matrices, operators)
\renewcommand{\AA}{{\mathbf A}}  %% Works but loses the Angstrom symbol
\newcommand{\BB}{{\mathbf B}}
\newcommand{\DD}{{\mathbf D}}
\newcommand{\EE}{{\mathbf E}}
\newcommand{\FF}{{\mathbf F}}
\newcommand{\GG}{{\mathbf G}}
\newcommand{\HH}{{\mathbf H}}
\newcommand{\II}{{\mathbf I}}
\newcommand{\JJ}{{\mathbf J}}
\newcommand{\KK}{{\mathbf K}}
\newcommand{\LL}{{\mathbf L}}
\newcommand{\MM}{{\mathbf M}}
\newcommand{\OO}{{\mathbf O}}
\newcommand{\PP}{{\mathbf P}}
\renewcommand{\SS}{{\mathbf S}}
\newcommand{\TT}{{\mathbf T}}
\newcommand{\UU}{{\mathbf U}}
\newcommand{\VV}{{\mathbf V}}
\newcommand{\WW}{{\mathbf W}}
\newcommand{\XX}{{\mathbf X}}
\newcommand{\YY}{{\mathbf Y}}

\renewcommand{\aa}{{\bf a}}    %% Redefines "å" to bold a
\newcommand{\bb}{{\bf b}}
\newcommand{\ff}{{\bf f}}
\newcommand{\jj}{{\bf j}}
\newcommand{\rr}{{\bf r}}
\newcommand{\pp}{{\bf p}}
\renewcommand{\ss}{{\bf s}}
\newcommand{\uu}{{\bf u}}
\newcommand{\vv}{{\bf v}}
\newcommand{\xx}{{\bf x}}
\newcommand{\yy}{{\bf y}}
\newcommand{\ee}{{\bf e}}

%% Bold Greek letters
\newcommand{\balpha}{\boldsymbol \alpha}
\newcommand{\bbeta}{\boldsymbol \beta}
\newcommand{\bomega}{\boldsymbol \omega}
\newcommand{\bOmega}{\boldsymbol \Omega}
\newcommand{\bgamma}{\boldsymbol \gamma}
\newcommand{\bvarphi}{\boldsymbol \varphi}
\newcommand{\blambda}{\boldsymbol \lambda}
\newcommand{\bsigma}{\boldsymbol \sigma}
\newcommand{\bpi}{\boldsymbol \pi}
\newcommand{\btau}{\boldsymbol \tau}
\newcommand{\bPhi}{\boldsymbol{\Phi}}

%%%%%%%%%%%%%%%%%%%%%%%%%%%%%%%%%%%%%%%%%%%%%%%%%%%%%%%%%%%%%%%%%%%%%%%%%%%%%%%
%%
%% End of Diffeology-macros.tex
%%
%%%%%%%%%%%%%%%%%%%%%%%%%%%%%%%%%%%%%%%%%%%%%%%%%%%%%%%%%%%%%%%%%%%%%%%%%%%%%%%

%%
%%  Date: 2026-02-21
%%
%%%%%%%%%%%%%%%%%%%%%%%%%%%%%%%%%%%%%%%%%%%%%%%%%%%%%%%%%%%%%%%%%%%%%%%%%%%%%%%

%%%%%%%%%%%%%%%%%%%%%%%%%%%%%%%%%%%%%%%%%%%%%%%%%%%%%%%%%%%%%%%%%%%%%%%%%%%%%%%
%%
%% MARK: Text Shortcuts and Hyphenation
%%
%%%%%%%%%%%%%%%%%%%%%%%%%%%%%%%%%%%%%%%%%%%%%%%%%%%%%%%%%%%%%%%%%%%%%%%%%%%%%%%

%% Implement text inside a display math block to make space around the text
\newcommand{\SpacedText}[1]{\ \text{#1}\ }
\newcommand{\QuadText}[1]{\quad\text{#1}\quad}

%% Hyphenation exceptions
\hyphenation{pa-ra-me-t-ri-za-tion}
\hyphenation{pa-ra-me-t-ri-za-tions}
\hyphenation{dif-feo-lo-gi-cal}
\hyphenation{dif-feo-lo-gy}

%% Latin shortcuts
\newcommand{\cf}{{\em cf.\/}}     %% Latin: confer, compare
\newcommand{\ie}{{\em i.e.\/}}    %% Latin: id est, that is
\newcommand{\resp}{{\em resp.\/}} %% Latin: respectively

%%%%%%%%%%%%%%%%%%%%%%%%%%%%%%%%%%%%%%%%%%%%%%%%%%%%%%%%%%%%%%%%%%%%%%%%%%%%%%%
%%
%% MARK: Number Sets
%%
%%%%%%%%%%%%%%%%%%%%%%%%%%%%%%%%%%%%%%%%%%%%%%%%%%%%%%%%%%%%%%%%%%%%%%%%%%%%%%%

%% Base command for number set formatting (French bold tradition)
\providecommand{\NumberSetBold}[1]{\mathbf{#1}}

%% ---------------------------------------------------------------------------
%% Level 1: Semantic names (for AI and future readers)
%% These explicitly state the mathematical nature of each set.
%% ---------------------------------------------------------------------------

%% \fieldR : The field of real numbers
\newcommand{\fieldR}{\NumberSetBold{R}}

%% \fieldC : The field of complex numbers
\newcommand{\fieldC}{\NumberSetBold{C}}

%% \fieldQ : The field of rational numbers
\newcommand{\fieldQ}{\NumberSetBold{Q}}

%% \ringZ : The ring of integers
\newcommand{\ringZ}{\NumberSetBold{Z}}

%% \naturalNumbers : The set of natural numbers (non-negative integers)
\newcommand{\naturalNumbers}{\NumberSetBold{N}}

%% ---------------------------------------------------------------------------
%% Level 2: Writing shortcuts (for the author's convenience)
%% These map your preferred concise notation to the semantic macros.
%% ---------------------------------------------------------------------------

%% \RR : Real numbers (shortcut for \fieldR)
\let\RR\fieldR

%% \CC : Complex numbers (shortcut for \fieldC)
\let\CC\fieldC

%% \QQ : Rational numbers (shortcut for \fieldQ)
\let\QQ\fieldQ

%% \ZZ : Integers (shortcut for \ringZ)
\let\ZZ\ringZ

%% \NN : Natural numbers (shortcut for \naturalNumbers)
\let\NN\naturalNumbers

%%%%%%%%%%%%%%%%%%%%%%%%%%%%%%%%%%%%%%%%%%%%%%%%%%%%%%%%%%%%%%%%%%%%%%%%%%%%%%%
%%
%% MARK: Categories and Morphisms
%%
%%%%%%%%%%%%%%%%%%%%%%%%%%%%%%%%%%%%%%%%%%%%%%%%%%%%%%%%%%%%%%%%%%%%%%%%%%%%%%%

%% \Category : Display the name of a category between braces
%% Usage: \Category{Top} produces {Top}
\newcommand{\Category}[1]{\{\text{#1}\}}

%% \Obj : The set of objects of a category
\newcommand{\Obj}{{\operatorname{Obj}}}

%% \Mor : The set of morphisms of a category
\newcommand{\Mor}{{\operatorname{Mor}}}

%% \Top : The set of topologies on a set (as in $\Top(\RR^n)$)
\newcommand{\Top}{\{\operatorname{Top}\}}

%% \Difflg : The set of all diffeologies on a set
\newcommand{\Difflg}{\operatorname{Diffeologies}}

%% \Hom : Homomorphisms (in any category)
\newcommand{\Hom}{\mathrm{Hom}}

%% \DHom : Smooth homomorphisms (in diffeological contexts)
\newcommand{\DHom}{\Hom^\infty}

%% \Isom : Isomorphisms
\newcommand{\Isom}{\mathrm{Isom}}

%% \DIsom : Smooth isomorphisms
\newcommand{\DIsom}{\Isom^\infty}

%%%%%%%%%%%%%%%%%%%%%%%%%%%%%%%%%%%%%%%%%%%%%%%%%%%%%%%%%%%%%%%%%%%%%%%%%%%%%%%
%%
%% MARK: Differential Calculus
%%
%%%%%%%%%%%%%%%%%%%%%%%%%%%%%%%%%%%%%%%%%%%%%%%%%%%%%%%%%%%%%%%%%%%%%%%%%%%%%%%

%% \Czero : Continuous maps
\newcommand{\Czero}{{\cC}^0}

%% \Ck : k-times continuously differentiable maps
%% Usage: \Ck{k} produces C^k
\newcommand{\Ck}[1]{{\cC}^{#1}}

%% \Cinfty : Infinitely differentiable (smooth) maps
\newcommand{\Cinfty}{{\cC}^\infty}

%% \Der : Tangent linear map (differential)
%% Type: \Der : (M → N) → M → T_x M → T_{f(x)} N
%% Usage: \Der(f)_x(v) or \Der(f)(x)(v)
%% Mathematical meaning: The derivative of smooth map f at point x,
%% applied to tangent vector v.
\newcommand{\Der}{\operatorname{D}}

%% \DForms : Differential forms on a diffeological space
%% Usage: \DForms^p(X) produces Ω^p(X)
\newcommand{\DForms}{\Omega}

%% Differential symbols for integration
\newcommand{\dt}{\mathrm{dt}}
\newcommand{\ds}{\mathrm{ds}}
\newcommand{\du}{\mathrm{du}}

%%%%%%%%%%%%%%%%%%%%%%%%%%%%%%%%%%%%%%%%%%%%%%%%%%%%%%%%%%%%%%%%%%%%%%%%%%%%%%%
%%
%% MARK: Chain-Homotopy Operator (Paths-Primitive)
%%
%%%%%%%%%%%%%%%%%%%%%%%%%%%%%%%%%%%%%%%%%%%%%%%%%%%%%%%%%%%%%%%%%%%%%%%%%%%%%%%

%% \CHO : Chain-Homotopy operator
\DeclareMathOperator{\CHO}{\boldsymbol{\mathsf{K}}}

%% \Komega : Paths-primitive of ω (1-form on Paths(X))
\newcommand{\Komega}{\CHO\omega}

%% \Kalpha : Paths-primitive of α
\newcommand{\Kalpha}{\CHO\alpha}

%% \Kbeta : Paths-primitive of β
\newcommand{\Kbeta}{\CHO\beta}

%%%%%%%%%%%%%%%%%%%%%%%%%%%%%%%%%%%%%%%%%%%%%%%%%%%%%%%%%%%%%%%%%%%%%%%%%%%%%%%
%%
%% MARK: Functions and Operators
%%
%%%%%%%%%%%%%%%%%%%%%%%%%%%%%%%%%%%%%%%%%%%%%%%%%%%%%%%%%%%%%%%%%%%%%%%%%%%%%%%

%% \Norm : Euclidean norm of a vector
\newcommand{\Norm}[1]{\Vert #1 \Vert}

%% \Modulus : Modulus of a complex number
\newcommand{\Modulus}[1]{\vert #1 \vert}

%% \Sign : Sign function
\newcommand{\Sign}{\operatorname{sgn}}

%% \det : Determinant
\renewcommand{\det}{\mathrm{det}}

%% \Sup : Supremum of a function
\newcommand{\Sup}{\operatorname{Sup}}

%% \Supp : Support of a function
%% Usage: \Supp(f) = closure of {x | f(x) ≠ 0}
\newcommand{\Supp}{\operatorname{Supp}}

%% \Id : Identity map
\newcommand{\Id}{{\mathbf{1}}}

%% \Eval : Evaluation map
%% Usage: \Eval(f,x) = f(x)
\newcommand{\Eval}{{\operatorname{ev}}}

%% \Rank : Rank of a linear map
\newcommand{\Rank}{\operatorname{rank}}

%% \Sq : Square norm operator
%% Usage: \Sq(x) = \Norm{x}^2
\newcommand{\Sq}{{\operatorname{sq}}}

%%%%%%%%%%%%%%%%%%%%%%%%%%%%%%%%%%%%%%%%%%%%%%%%%%%%%%%%%%%%%%%%%%%%%%%%%%%%%%%
%%
%% MARK: Groups and Lie Theory
%%
%%%%%%%%%%%%%%%%%%%%%%%%%%%%%%%%%%%%%%%%%%%%%%%%%%%%%%%%%%%%%%%%%%%%%%%%%%%%%%%

%% Linear and matrix groups
\newcommand{\GL}{\mathrm{GL}}     %% General linear group
\newcommand{\SL}{\mathrm{SL}}     %% Special linear group
\newcommand{\SO}{\mathrm{SO}}     %% Special orthogonal group
\renewcommand{\O}{\mathrm{O}}     %% Orthogonal group
\newcommand{\Sp}{{\rm Sp}}        %% Symplectic group
\newcommand{\U}{\mathrm{U}}       %% Unitary group
\newcommand{\SU}{\mathrm{SU}}     %% Special unitary group
\newcommand{\Perm}{\mathrm{Perm}} %% Permutation group

%% Transformation groups
\newcommand{\Aff}{\operatorname{Aff}}   %% Affine group
\newcommand{\Aut}{\operatorname{Aut}}   %% Automorphism group
\newcommand{\Diff}{\operatorname{Diff}} %% Diffeomorphism group
\newcommand{\Ham}{\operatorname{Ham}}   %% Hamiltonian diffeomorphisms

%% Group operations
\newcommand{\Ad}{\mathrm{Ad}}   %% Adjoint action: Ad(g)(k) = gkg^{-1}
\newcommand{\Lm}{\operatorname{L}}   %% Left multiplication: L(g)(k) = gk
\newcommand{\Rm}{\operatorname{R}}   %% Right multiplication: R(g)(k) = kg

%% \OrbitMap : Orbit map of a group action
%% Usage: \OrbitMap{x}(g) = g(x)
\newcommand{\OrbitMap}[1]{\hat{#1}}

%% \Stab : Stabilizer of a point under a group action
\newcommand{\Stab}{\mathrm{St}}

%% \LieAlgebra : Lie algebra of a Lie group
%% Usage: \LieAlgebra{G} produces \mathfrak{g}
\newcommand{\LieAlgebra}[1]{\mathfrak{\MakeLowercase{#1}}}

%% \DLie : Lie derivative (pounds sign)
\newcommand{\DLie}{\mbox{\pounds}}

%%%%%%%%%%%%%%%%%%%%%%%%%%%%%%%%%%%%%%%%%%%%%%%%%%%%%%%%%%%%%%%%%%%%%%%%%%%%%%%
%%
%% MARK: Tangent Space and Vector Fields
%%
%%%%%%%%%%%%%%%%%%%%%%%%%%%%%%%%%%%%%%%%%%%%%%%%%%%%%%%%%%%%%%%%%%%%%%%%%%%%%%%

%% \Tgt : Tangent space or tangent bundle
%% Usage: \Tgt_x(X) for tangent space at point x
%%        \Tgt X for tangent bundle of X
\newcommand{\Tgt}{T}

%% \TgtBundle : Explicit tangent bundle notation (verbose)
\newcommand{\TgtBundle}{T}

%% \UnitBundle : Unit tangent bundle
\newcommand{\UnitBundle}{U}

%% \VectorField : Space of vector fields on a manifold
\newcommand{\VectorField}{\operatorname{Vect}}

%% \grad : Gradient vector field
\newcommand{\grad}{\operatorname{grad}}

%%%%%%%%%%%%%%%%%%%%%%%%%%%%%%%%%%%%%%%%%%%%%%%%%%%%%%%%%%%%%%%%%%%%%%%%%%%%%%%
%%
%% MARK: Paths, Loops, and Homotopy
%%
%%%%%%%%%%%%%%%%%%%%%%%%%%%%%%%%%%%%%%%%%%%%%%%%%%%%%%%%%%%%%%%%%%%%%%%%%%%%%%%

%% \Paths : Space of smooth paths in a diffeological space
\newcommand{\Paths}{\operatorname{Paths}}

%% \Loops : Space of smooth loops (paths with coincident endpoints)
\newcommand{\Loops}{\operatorname{Loops}}

%% \stPaths : Space of stationary paths
\newcommand{\stPaths}{\operatorname{stPaths}}

%% \stLoops : Space of stationary loops
\newcommand{\stLoops}{\operatorname{stLoops}}

%% \0 : Origin (starting point) of a path
\newcommand{\0}{\hat 0}

%% \1 : End (terminal point) of a path
\newcommand{\1}{\hat 1}

%% \Ends : Pair of endpoints of a path
\def\Ends{\operatorname{ends}}

%% \Rev : Reverse of a path
%% Usage: \Rev(\gamma)(t) = \gamma(1-t)
\def\Rev{\operatorname{rev}}

%% \Comp : Connected component containing a point
\newcommand{\Comp}{\operatorname{comp}}

%%%%%%%%%%%%%%%%%%%%%%%%%%%%%%%%%%%%%%%%%%%%%%%%%%%%%%%%%%%%%%%%%%%%%%%%%%%%%%%
%%
%% MARK: Homology and Cohomology
%%
%%%%%%%%%%%%%%%%%%%%%%%%%%%%%%%%%%%%%%%%%%%%%%%%%%%%%%%%%%%%%%%%%%%%%%%%%%%%%%%

%% \Boundary : Boundary operator in homology
\DeclareMathOperator{\Boundary}{\partial}

%% \cub : The p-dimensional cube I^p (I = [0,1])
\newcommand{\cub}[1]{\operatorname{I}^{#1}}

%% \DCubes : Space of smooth cubes in X
\newcommand{\DCubes}[1]{\operatorname{Cub}_{#1}}

%% \DChains : Smooth cubic chains
\newcommand{\DChains}[1]{\operatorname{C}_{#1}}

%% \DCochains : Smooth cubic cochains
\newcommand{\DCochains}[1]{\operatorname{C}^{#1}}

%% Reduced (quotiented by degenerates) chain groups
\newcommand{\QDChains}[1]{\operatorname{\textbf{\textsf C}}_{#1}}
\newcommand{\QDCochains}[1]{\operatorname{\textbf{\textsf C}}^{#1}}

%% Cycles, boundaries, and homology
\newcommand{\QDZ}{\operatorname{\textbf{\textsf Z}}}   %% Cycles/cocycles
\newcommand{\QDB}{\operatorname{\textbf{\textsf B}}}   %% Boundaries/coboundaries
\newcommand{\QDH}{\operatorname{\textbf{\textsf H}}}   %% Homology/cohomology groups

%% Generic homology notation
\newcommand{\ZG}{Z}   %% Cycles
\newcommand{\BG}{B}   %% Boundaries
\newcommand{\HG}{H}   %% Homology groups

%% De Rham cohomology
\newcommand{\ZDR}{Z_\mathrm{dR}}   %% De Rham cocycles
\newcommand{\BDR}{B_\mathrm{dR}}   %% De Rham coboundaries
\newcommand{\HDR}{H_\mathrm{dR}}   %% De Rham cohomology

%% \hdR : De Rham homomorphism
\newcommand{\hdR}{h_{\mathrm{dR}}}

%%%%%%%%%%%%%%%%%%%%%%%%%%%%%%%%%%%%%%%%%%%%%%%%%%%%%%%%%%%%%%%%%%%%%%%%%%%%%%%
%%
%% MARK: Diffeology Specific
%%
%%%%%%%%%%%%%%%%%%%%%%%%%%%%%%%%%%%%%%%%%%%%%%%%%%%%%%%%%%%%%%%%%%%%%%%%%%%%%%%

%% \Domain : Domain of a parametrization
\newcommand{\Domain}{\mathrm{dom}}

%% \Domains : Set of domains in Euclidean space
\newcommand{\Domains}{\mathrm{Domains}}

%% \Params : Set of parametrizations in a set
\newcommand{\Params}{\mathrm{Params}}

%% \Plots : Set of plots in a diffeological space (the diffeology)
\newcommand{\Plots}{\mathrm{Plots}}

%%%%%%%%%%%%%%%%%%%%%%%%%%%%%%%%%%%%%%%%%%%%%%%%%%%%%%%%%%%%%%%%%%%%%%%%%%%%%%%
%%
%% MARK: Quotients and Projections
%%
%%%%%%%%%%%%%%%%%%%%%%%%%%%%%%%%%%%%%%%%%%%%%%%%%%%%%%%%%%%%%%%%%%%%%%%%%%%%%%%

%% \Class : The equivalence class of an element
%% Usage: \Class(x) = { y | y ~ x }
\newcommand{\Class}{\operatorname{class}}

%% \Cl : Short notation for the class of an element
%% Usage: \Cl[x] = \Class(x)
\newcommand{\Cl}[1]{\left[ #1 \right]}

%% \Proj : The canonical projection map from a set to its quotient
%% Usage: \Proj : X → X/∼
%% For a specific quotient: \Proj_{\Gamma} : ℝ → ℝ/Γ
\newcommand{\Proj}{\operatorname{pr}}

%%%%%%%%%%%%%%%%%%%%%%%%%%%%%%%%%%%%%%%%%%%%%%%%%%%%%%%%%%%%%%%%%%%%%%%%%%%%%%%
%%
%% MARK: Covering Spaces
%%
%%%%%%%%%%%%%%%%%%%%%%%%%%%%%%%%%%%%%%%%%%%%%%%%%%%%%%%%%%%%%%%%%%%%%%%%%%%%%%%

%% \UniversalCover : Universal covering space (maximal simply connected cover)
%% Uses the stretching tilde from AMS.
%% Usage: \UniversalCover{\cT_{\alpha}} produces \widetilde{\cT_{\alpha}}
\newcommand{\UniversalCover}[1]{\widetilde{#1}}

%% \Cover : Any covering space (not necessarily universal)
%% Uses the hat from AMS. This is the general covering projection.
%% Usage: \Cover{G} produces \widehat{G}
\newcommand{\Cover}[1]{\widehat{#1}}

%%%%%%%%%%%%%%%%%%%%%%%%%%%%%%%%%%%%%%%%%%%%%%%%%%%%%%%%%%%%%%%%%%%%%%%%%%%%%%%
%%
%% MARK: Homotopy and Fundamental Group
%%
%%%%%%%%%%%%%%%%%%%%%%%%%%%%%%%%%%%%%%%%%%%%%%%%%%%%%%%%%%%%%%%%%%%%%%%%%%%%%%%

%% \FundamentalGroup : The fundamental group (first homotopy group)
%% Usage: \FundamentalGroup{X} produces π₁(X)
%% For a pointed space, use \FundamentalGroup{X,x₀}
\newcommand{\FundamentalGroup}[1]{\pi_1(#1)}

%% \piX : Alternative semantic name (shorter, still clear)
%% Usage: \piX{X} produces π₁(X)
\newcommand{\piX}[1]{\pi_1(#1)}

%%%%%%%%%%%%%%%%%%%%%%%%%%%%%%%%%%%%%%%%%%%%%%%%%%%%%%%%%%%%%%%%%%%%%%%%%%%%%%%
%%
%% MARK: Periods and Torus
%%
%%%%%%%%%%%%%%%%%%%%%%%%%%%%%%%%%%%%%%%%%%%%%%%%%%%%%%%%%%%%%%%%%%%%%%%%%%%%%%%

%% \Periods : Group of periods
\newcommand{\Periods}{\operatorname{P}}

%% \cTorus : Calligraphic T for distinctive infinite torus
\newcommand{\cTorus}{\cT}

%% \Torus : Standard torus R/Z
\newcommand{\Torus}{\mathrm{T}}

%% \Sph : n-sphere
\newcommand{\Sph}{\mathrm{S}}

%%%%%%%%%%%%%%%%%%%%%%%%%%%%%%%%%%%%%%%%%%%%%%%%%%%%%%%%%%%%%%%%%%%%%%%%%%%%%%%
%%
%% MARK: Linear Algebra and Vector Spaces
%%
%%%%%%%%%%%%%%%%%%%%%%%%%%%%%%%%%%%%%%%%%%%%%%%%%%%%%%%%%%%%%%%%%%%%%%%%%%%%%%%

%% \Lin : Space of linear maps between vector spaces
\newcommand{\Lin}{\operatorname{L}}

%% \DLin : Smooth linear maps between diffeological vector spaces
\newcommand{\DLin}{\Lin^\infty}

%% \Vector : Column vector with parentheses
%% Usage: \Vector{a \\ b \\ c}
\newcommand{\Vector}[1]{\begin{pmatrix}#1\end{pmatrix}}

%% \SqVect : Column vector with square brackets
\newcommand{\SqVect}[1]{\left[\begin{matrix}#1\end{matrix}\right]}

%% \Matrix : Set of matrices
\DeclareMathOperator{\Matrix}{Mat}

%% \ker : Kernel of a linear map
\renewcommand{\ker}{\mathrm{ker}}

%% \coker : Cokernel of a linear map
\newcommand{\coker}{\operatorname{coker}}

%% \J : The π/2 rotation matrix (complex structure)
\newcommand{\J}{\operatorname{J}}

%% \vecspan : Linear span of vectors
\newcommand{\vecspan}[1]{{\rm span}(#1)}

%%%%%%%%%%%%%%%%%%%%%%%%%%%%%%%%%%%%%%%%%%%%%%%%%%%%%%%%%%%%%%%%%%%%%%%%%%%%%%%
%%
%% MARK: Fonts (Semantic Variants)
%%
%% These provide semantic aliases for font variants.
%% The visual appearance is handled by Display-macros.tex
%%
%%%%%%%%%%%%%%%%%%%%%%%%%%%%%%%%%%%%%%%%%%%%%%%%%%%%%%%%%%%%%%%%%%%%%%%%%%%%%%%

%% Calligraphic letters (semantic: used for categories, sheaves, etc.)
\newcommand{\cA}{\mathcal{A}}
\newcommand{\cB}{\mathcal{B}}
\newcommand{\cC}{\mathcal{C}}
\newcommand{\cD}{\mathcal{D}}
\newcommand{\cE}{\mathcal{E}}
\newcommand{\cF}{\mathcal{F}}
\newcommand{\cG}{\mathcal{G}}
\newcommand{\cH}{\mathcal{H}}
\newcommand{\cI}{\mathcal{I}}
\newcommand{\cJ}{\mathcal{J}}
\newcommand{\cK}{\mathcal{K}}
\newcommand{\cL}{\mathcal{L}}
\newcommand{\cM}{\mathcal{M}}
\newcommand{\cN}{\mathcal{N}}
\newcommand{\cO}{\mathcal{O}}
\newcommand{\cP}{\mathcal{P}}
\newcommand{\cQ}{\mathcal{Q}}
\newcommand{\cR}{\mathcal{R}}
\newcommand{\cS}{\mathcal{S}}
\newcommand{\cT}{\mathcal{T}}
\newcommand{\cU}{\mathcal{U}}
\newcommand{\cV}{\mathcal{V}}
\newcommand{\cW}{\mathcal{W}}
\newcommand{\cX}{\mathcal{X}}
\newcommand{\cY}{\mathcal{Y}}
\newcommand{\cZ}{\mathcal{Z}}

%% Fraktur letters (semantic: used for Lie algebras, ideals, etc.)
\newcommand{\fA}{{\mathfrak A}}
\newcommand{\fB}{{\mathfrak B}}
\newcommand{\fC}{{\mathfrak C}}
\newcommand{\fD}{{\mathfrak D}}
\newcommand{\fF}{{\mathfrak F}}
\newcommand{\fG}{{\mathfrak G}}
\newcommand{\fK}{{\mathfrak K}}
\newcommand{\fI}{{\mathfrak I}}
\newcommand{\fL}{{\mathfrak L}}
\newcommand{\fP}{{\mathfrak P}}
\newcommand{\fR}{{\mathfrak R}}
\newcommand{\fS}{{\mathfrak S}}
\newcommand{\fU}{{\mathfrak U}}
\newcommand{\fX}{{\mathfrak X}}
\newcommand{\fY}{{\mathfrak Y}}
\newcommand{\fZ}{{\mathfrak Z}}
\newcommand{\fu}{{\mathfrak u}}

%% Bold letters (semantic: used for vectors, matrices, operators)
\renewcommand{\AA}{{\mathbf A}}  %% Works but loses the Angstrom symbol
\newcommand{\BB}{{\mathbf B}}
\newcommand{\DD}{{\mathbf D}}
\newcommand{\EE}{{\mathbf E}}
\newcommand{\FF}{{\mathbf F}}
\newcommand{\GG}{{\mathbf G}}
\newcommand{\HH}{{\mathbf H}}
\newcommand{\II}{{\mathbf I}}
\newcommand{\JJ}{{\mathbf J}}
\newcommand{\KK}{{\mathbf K}}
\newcommand{\LL}{{\mathbf L}}
\newcommand{\MM}{{\mathbf M}}
\newcommand{\OO}{{\mathbf O}}
\newcommand{\PP}{{\mathbf P}}
\renewcommand{\SS}{{\mathbf S}}
\newcommand{\TT}{{\mathbf T}}
\newcommand{\UU}{{\mathbf U}}
\newcommand{\VV}{{\mathbf V}}
\newcommand{\WW}{{\mathbf W}}
\newcommand{\XX}{{\mathbf X}}
\newcommand{\YY}{{\mathbf Y}}

\renewcommand{\aa}{{\bf a}}    %% Redefines "å" to bold a
\newcommand{\bb}{{\bf b}}
\newcommand{\ff}{{\bf f}}
\newcommand{\jj}{{\bf j}}
\newcommand{\rr}{{\bf r}}
\newcommand{\pp}{{\bf p}}
\renewcommand{\ss}{{\bf s}}
\newcommand{\uu}{{\bf u}}
\newcommand{\vv}{{\bf v}}
\newcommand{\xx}{{\bf x}}
\newcommand{\yy}{{\bf y}}
\newcommand{\ee}{{\bf e}}

%% Bold Greek letters
\newcommand{\balpha}{\boldsymbol \alpha}
\newcommand{\bbeta}{\boldsymbol \beta}
\newcommand{\bomega}{\boldsymbol \omega}
\newcommand{\bOmega}{\boldsymbol \Omega}
\newcommand{\bgamma}{\boldsymbol \gamma}
\newcommand{\bvarphi}{\boldsymbol \varphi}
\newcommand{\blambda}{\boldsymbol \lambda}
\newcommand{\bsigma}{\boldsymbol \sigma}
\newcommand{\bpi}{\boldsymbol \pi}
\newcommand{\btau}{\boldsymbol \tau}
\newcommand{\bPhi}{\boldsymbol{\Phi}}

%%%%%%%%%%%%%%%%%%%%%%%%%%%%%%%%%%%%%%%%%%%%%%%%%%%%%%%%%%%%%%%%%%%%%%%%%%%%%%%
%%
%% End of Diffeology-macros.tex
%%
%%%%%%%%%%%%%%%%%%%%%%%%%%%%%%%%%%%%%%%%%%%%%%%%%%%%%%%%%%%%%%%%%%%%%%%%%%%%%%%

%%
%%  Date: 2026-02-21
%%
%%%%%%%%%%%%%%%%%%%%%%%%%%%%%%%%%%%%%%%%%%%%%%%%%%%%%%%%%%%%%%%%%%%%%%%%%%%%%%%

%%%%%%%%%%%%%%%%%%%%%%%%%%%%%%%%%%%%%%%%%%%%%%%%%%%%%%%%%%%%%%%%%%%%%%%%%%%%%%%
%%
%% MARK: Text Shortcuts and Hyphenation
%%
%%%%%%%%%%%%%%%%%%%%%%%%%%%%%%%%%%%%%%%%%%%%%%%%%%%%%%%%%%%%%%%%%%%%%%%%%%%%%%%

%% Implement text inside a display math block to make space around the text
\newcommand{\SpacedText}[1]{\ \text{#1}\ }
\newcommand{\QuadText}[1]{\quad\text{#1}\quad}

%% Hyphenation exceptions
\hyphenation{pa-ra-me-t-ri-za-tion}
\hyphenation{pa-ra-me-t-ri-za-tions}
\hyphenation{dif-feo-lo-gi-cal}
\hyphenation{dif-feo-lo-gy}

%% Latin shortcuts
\newcommand{\cf}{{\em cf.\/}}     %% Latin: confer, compare
\newcommand{\ie}{{\em i.e.\/}}    %% Latin: id est, that is
\newcommand{\resp}{{\em resp.\/}} %% Latin: respectively

%%%%%%%%%%%%%%%%%%%%%%%%%%%%%%%%%%%%%%%%%%%%%%%%%%%%%%%%%%%%%%%%%%%%%%%%%%%%%%%
%%
%% MARK: Number Sets
%%
%%%%%%%%%%%%%%%%%%%%%%%%%%%%%%%%%%%%%%%%%%%%%%%%%%%%%%%%%%%%%%%%%%%%%%%%%%%%%%%

%% Base command for number set formatting (French bold tradition)
\providecommand{\NumberSetBold}[1]{\mathbf{#1}}

%% ---------------------------------------------------------------------------
%% Level 1: Semantic names (for AI and future readers)
%% These explicitly state the mathematical nature of each set.
%% ---------------------------------------------------------------------------

%% \fieldR : The field of real numbers
\newcommand{\fieldR}{\NumberSetBold{R}}

%% \fieldC : The field of complex numbers
\newcommand{\fieldC}{\NumberSetBold{C}}

%% \fieldQ : The field of rational numbers
\newcommand{\fieldQ}{\NumberSetBold{Q}}

%% \ringZ : The ring of integers
\newcommand{\ringZ}{\NumberSetBold{Z}}

%% \naturalNumbers : The set of natural numbers (non-negative integers)
\newcommand{\naturalNumbers}{\NumberSetBold{N}}

%% ---------------------------------------------------------------------------
%% Level 2: Writing shortcuts (for the author's convenience)
%% These map your preferred concise notation to the semantic macros.
%% ---------------------------------------------------------------------------

%% \RR : Real numbers (shortcut for \fieldR)
\let\RR\fieldR

%% \CC : Complex numbers (shortcut for \fieldC)
\let\CC\fieldC

%% \QQ : Rational numbers (shortcut for \fieldQ)
\let\QQ\fieldQ

%% \ZZ : Integers (shortcut for \ringZ)
\let\ZZ\ringZ

%% \NN : Natural numbers (shortcut for \naturalNumbers)
\let\NN\naturalNumbers

%%%%%%%%%%%%%%%%%%%%%%%%%%%%%%%%%%%%%%%%%%%%%%%%%%%%%%%%%%%%%%%%%%%%%%%%%%%%%%%
%%
%% MARK: Categories and Morphisms
%%
%%%%%%%%%%%%%%%%%%%%%%%%%%%%%%%%%%%%%%%%%%%%%%%%%%%%%%%%%%%%%%%%%%%%%%%%%%%%%%%

%% \Category : Display the name of a category between braces
%% Usage: \Category{Top} produces {Top}
\newcommand{\Category}[1]{\{\text{#1}\}}

%% \Obj : The set of objects of a category
\newcommand{\Obj}{{\operatorname{Obj}}}

%% \Mor : The set of morphisms of a category
\newcommand{\Mor}{{\operatorname{Mor}}}

%% \Top : The set of topologies on a set (as in $\Top(\RR^n)$)
\newcommand{\Top}{\{\operatorname{Top}\}}

%% \Difflg : The set of all diffeologies on a set
\newcommand{\Difflg}{\operatorname{Diffeologies}}

%% \Hom : Homomorphisms (in any category)
\newcommand{\Hom}{\mathrm{Hom}}

%% \DHom : Smooth homomorphisms (in diffeological contexts)
\newcommand{\DHom}{\Hom^\infty}

%% \Isom : Isomorphisms
\newcommand{\Isom}{\mathrm{Isom}}

%% \DIsom : Smooth isomorphisms
\newcommand{\DIsom}{\Isom^\infty}

%%%%%%%%%%%%%%%%%%%%%%%%%%%%%%%%%%%%%%%%%%%%%%%%%%%%%%%%%%%%%%%%%%%%%%%%%%%%%%%
%%
%% MARK: Differential Calculus
%%
%%%%%%%%%%%%%%%%%%%%%%%%%%%%%%%%%%%%%%%%%%%%%%%%%%%%%%%%%%%%%%%%%%%%%%%%%%%%%%%

%% \Czero : Continuous maps
\newcommand{\Czero}{{\cC}^0}

%% \Ck : k-times continuously differentiable maps
%% Usage: \Ck{k} produces C^k
\newcommand{\Ck}[1]{{\cC}^{#1}}

%% \Cinfty : Infinitely differentiable (smooth) maps
\newcommand{\Cinfty}{{\cC}^\infty}

%% \Der : Tangent linear map (differential)
%% Type: \Der : (M → N) → M → T_x M → T_{f(x)} N
%% Usage: \Der(f)_x(v) or \Der(f)(x)(v)
%% Mathematical meaning: The derivative of smooth map f at point x,
%% applied to tangent vector v.
\newcommand{\Der}{\operatorname{D}}

%% \DForms : Differential forms on a diffeological space
%% Usage: \DForms^p(X) produces Ω^p(X)
\newcommand{\DForms}{\Omega}

%% Differential symbols for integration
\newcommand{\dt}{\mathrm{dt}}
\newcommand{\ds}{\mathrm{ds}}
\newcommand{\du}{\mathrm{du}}

%%%%%%%%%%%%%%%%%%%%%%%%%%%%%%%%%%%%%%%%%%%%%%%%%%%%%%%%%%%%%%%%%%%%%%%%%%%%%%%
%%
%% MARK: Chain-Homotopy Operator (Paths-Primitive)
%%
%%%%%%%%%%%%%%%%%%%%%%%%%%%%%%%%%%%%%%%%%%%%%%%%%%%%%%%%%%%%%%%%%%%%%%%%%%%%%%%

%% \CHO : Chain-Homotopy operator
\DeclareMathOperator{\CHO}{\boldsymbol{\mathsf{K}}}

%% \Komega : Paths-primitive of ω (1-form on Paths(X))
\newcommand{\Komega}{\CHO\omega}

%% \Kalpha : Paths-primitive of α
\newcommand{\Kalpha}{\CHO\alpha}

%% \Kbeta : Paths-primitive of β
\newcommand{\Kbeta}{\CHO\beta}

%%%%%%%%%%%%%%%%%%%%%%%%%%%%%%%%%%%%%%%%%%%%%%%%%%%%%%%%%%%%%%%%%%%%%%%%%%%%%%%
%%
%% MARK: Functions and Operators
%%
%%%%%%%%%%%%%%%%%%%%%%%%%%%%%%%%%%%%%%%%%%%%%%%%%%%%%%%%%%%%%%%%%%%%%%%%%%%%%%%

%% \Norm : Euclidean norm of a vector
\newcommand{\Norm}[1]{\Vert #1 \Vert}

%% \Modulus : Modulus of a complex number
\newcommand{\Modulus}[1]{\vert #1 \vert}

%% \Sign : Sign function
\newcommand{\Sign}{\operatorname{sgn}}

%% \det : Determinant
\renewcommand{\det}{\mathrm{det}}

%% \Sup : Supremum of a function
\newcommand{\Sup}{\operatorname{Sup}}

%% \Supp : Support of a function
%% Usage: \Supp(f) = closure of {x | f(x) ≠ 0}
\newcommand{\Supp}{\operatorname{Supp}}

%% \Id : Identity map
\newcommand{\Id}{{\mathbf{1}}}

%% \Eval : Evaluation map
%% Usage: \Eval(f,x) = f(x)
\newcommand{\Eval}{{\operatorname{ev}}}

%% \Rank : Rank of a linear map
\newcommand{\Rank}{\operatorname{rank}}

%% \Sq : Square norm operator
%% Usage: \Sq(x) = \Norm{x}^2
\newcommand{\Sq}{{\operatorname{sq}}}

%%%%%%%%%%%%%%%%%%%%%%%%%%%%%%%%%%%%%%%%%%%%%%%%%%%%%%%%%%%%%%%%%%%%%%%%%%%%%%%
%%
%% MARK: Groups and Lie Theory
%%
%%%%%%%%%%%%%%%%%%%%%%%%%%%%%%%%%%%%%%%%%%%%%%%%%%%%%%%%%%%%%%%%%%%%%%%%%%%%%%%

%% Linear and matrix groups
\newcommand{\GL}{\mathrm{GL}}     %% General linear group
\newcommand{\SL}{\mathrm{SL}}     %% Special linear group
\newcommand{\SO}{\mathrm{SO}}     %% Special orthogonal group
\renewcommand{\O}{\mathrm{O}}     %% Orthogonal group
\newcommand{\Sp}{{\rm Sp}}        %% Symplectic group
\newcommand{\U}{\mathrm{U}}       %% Unitary group
\newcommand{\SU}{\mathrm{SU}}     %% Special unitary group
\newcommand{\Perm}{\mathrm{Perm}} %% Permutation group

%% Transformation groups
\newcommand{\Aff}{\operatorname{Aff}}   %% Affine group
\newcommand{\Aut}{\operatorname{Aut}}   %% Automorphism group
\newcommand{\Diff}{\operatorname{Diff}} %% Diffeomorphism group
\newcommand{\Ham}{\operatorname{Ham}}   %% Hamiltonian diffeomorphisms

%% Group operations
\newcommand{\Ad}{\mathrm{Ad}}   %% Adjoint action: Ad(g)(k) = gkg^{-1}
\newcommand{\Lm}{\operatorname{L}}   %% Left multiplication: L(g)(k) = gk
\newcommand{\Rm}{\operatorname{R}}   %% Right multiplication: R(g)(k) = kg

%% \OrbitMap : Orbit map of a group action
%% Usage: \OrbitMap{x}(g) = g(x)
\newcommand{\OrbitMap}[1]{\hat{#1}}

%% \Stab : Stabilizer of a point under a group action
\newcommand{\Stab}{\mathrm{St}}

%% \LieAlgebra : Lie algebra of a Lie group
%% Usage: \LieAlgebra{G} produces \mathfrak{g}
\newcommand{\LieAlgebra}[1]{\mathfrak{\MakeLowercase{#1}}}

%% \DLie : Lie derivative (pounds sign)
\newcommand{\DLie}{\mbox{\pounds}}

%%%%%%%%%%%%%%%%%%%%%%%%%%%%%%%%%%%%%%%%%%%%%%%%%%%%%%%%%%%%%%%%%%%%%%%%%%%%%%%
%%
%% MARK: Tangent Space and Vector Fields
%%
%%%%%%%%%%%%%%%%%%%%%%%%%%%%%%%%%%%%%%%%%%%%%%%%%%%%%%%%%%%%%%%%%%%%%%%%%%%%%%%

%% \Tgt : Tangent space or tangent bundle
%% Usage: \Tgt_x(X) for tangent space at point x
%%        \Tgt X for tangent bundle of X
\newcommand{\Tgt}{T}

%% \TgtBundle : Explicit tangent bundle notation (verbose)
\newcommand{\TgtBundle}{T}

%% \UnitBundle : Unit tangent bundle
\newcommand{\UnitBundle}{U}

%% \VectorField : Space of vector fields on a manifold
\newcommand{\VectorField}{\operatorname{Vect}}

%% \grad : Gradient vector field
\newcommand{\grad}{\operatorname{grad}}

%%%%%%%%%%%%%%%%%%%%%%%%%%%%%%%%%%%%%%%%%%%%%%%%%%%%%%%%%%%%%%%%%%%%%%%%%%%%%%%
%%
%% MARK: Paths, Loops, and Homotopy
%%
%%%%%%%%%%%%%%%%%%%%%%%%%%%%%%%%%%%%%%%%%%%%%%%%%%%%%%%%%%%%%%%%%%%%%%%%%%%%%%%

%% \Paths : Space of smooth paths in a diffeological space
\newcommand{\Paths}{\operatorname{Paths}}

%% \Loops : Space of smooth loops (paths with coincident endpoints)
\newcommand{\Loops}{\operatorname{Loops}}

%% \stPaths : Space of stationary paths
\newcommand{\stPaths}{\operatorname{stPaths}}

%% \stLoops : Space of stationary loops
\newcommand{\stLoops}{\operatorname{stLoops}}

%% \0 : Origin (starting point) of a path
\newcommand{\0}{\hat 0}

%% \1 : End (terminal point) of a path
\newcommand{\1}{\hat 1}

%% \Ends : Pair of endpoints of a path
\def\Ends{\operatorname{ends}}

%% \Rev : Reverse of a path
%% Usage: \Rev(\gamma)(t) = \gamma(1-t)
\def\Rev{\operatorname{rev}}

%% \Comp : Connected component containing a point
\newcommand{\Comp}{\operatorname{comp}}

%%%%%%%%%%%%%%%%%%%%%%%%%%%%%%%%%%%%%%%%%%%%%%%%%%%%%%%%%%%%%%%%%%%%%%%%%%%%%%%
%%
%% MARK: Homology and Cohomology
%%
%%%%%%%%%%%%%%%%%%%%%%%%%%%%%%%%%%%%%%%%%%%%%%%%%%%%%%%%%%%%%%%%%%%%%%%%%%%%%%%

%% \Boundary : Boundary operator in homology
\DeclareMathOperator{\Boundary}{\partial}

%% \cub : The p-dimensional cube I^p (I = [0,1])
\newcommand{\cub}[1]{\operatorname{I}^{#1}}

%% \DCubes : Space of smooth cubes in X
\newcommand{\DCubes}[1]{\operatorname{Cub}_{#1}}

%% \DChains : Smooth cubic chains
\newcommand{\DChains}[1]{\operatorname{C}_{#1}}

%% \DCochains : Smooth cubic cochains
\newcommand{\DCochains}[1]{\operatorname{C}^{#1}}

%% Reduced (quotiented by degenerates) chain groups
\newcommand{\QDChains}[1]{\operatorname{\textbf{\textsf C}}_{#1}}
\newcommand{\QDCochains}[1]{\operatorname{\textbf{\textsf C}}^{#1}}

%% Cycles, boundaries, and homology
\newcommand{\QDZ}{\operatorname{\textbf{\textsf Z}}}   %% Cycles/cocycles
\newcommand{\QDB}{\operatorname{\textbf{\textsf B}}}   %% Boundaries/coboundaries
\newcommand{\QDH}{\operatorname{\textbf{\textsf H}}}   %% Homology/cohomology groups

%% Generic homology notation
\newcommand{\ZG}{Z}   %% Cycles
\newcommand{\BG}{B}   %% Boundaries
\newcommand{\HG}{H}   %% Homology groups

%% De Rham cohomology
\newcommand{\ZDR}{Z_\mathrm{dR}}   %% De Rham cocycles
\newcommand{\BDR}{B_\mathrm{dR}}   %% De Rham coboundaries
\newcommand{\HDR}{H_\mathrm{dR}}   %% De Rham cohomology

%% \hdR : De Rham homomorphism
\newcommand{\hdR}{h_{\mathrm{dR}}}

%%%%%%%%%%%%%%%%%%%%%%%%%%%%%%%%%%%%%%%%%%%%%%%%%%%%%%%%%%%%%%%%%%%%%%%%%%%%%%%
%%
%% MARK: Diffeology Specific
%%
%%%%%%%%%%%%%%%%%%%%%%%%%%%%%%%%%%%%%%%%%%%%%%%%%%%%%%%%%%%%%%%%%%%%%%%%%%%%%%%

%% \Domain : Domain of a parametrization
\newcommand{\Domain}{\mathrm{dom}}

%% \Domains : Set of domains in Euclidean space
\newcommand{\Domains}{\mathrm{Domains}}

%% \Params : Set of parametrizations in a set
\newcommand{\Params}{\mathrm{Params}}

%% \Plots : Set of plots in a diffeological space (the diffeology)
\newcommand{\Plots}{\mathrm{Plots}}

%%%%%%%%%%%%%%%%%%%%%%%%%%%%%%%%%%%%%%%%%%%%%%%%%%%%%%%%%%%%%%%%%%%%%%%%%%%%%%%
%%
%% MARK: Quotients and Projections
%%
%%%%%%%%%%%%%%%%%%%%%%%%%%%%%%%%%%%%%%%%%%%%%%%%%%%%%%%%%%%%%%%%%%%%%%%%%%%%%%%

%% \Class : The equivalence class of an element
%% Usage: \Class(x) = { y | y ~ x }
\newcommand{\Class}{\operatorname{class}}

%% \Cl : Short notation for the class of an element
%% Usage: \Cl[x] = \Class(x)
\newcommand{\Cl}[1]{\left[ #1 \right]}

%% \Proj : The canonical projection map from a set to its quotient
%% Usage: \Proj : X → X/∼
%% For a specific quotient: \Proj_{\Gamma} : ℝ → ℝ/Γ
\newcommand{\Proj}{\operatorname{pr}}

%%%%%%%%%%%%%%%%%%%%%%%%%%%%%%%%%%%%%%%%%%%%%%%%%%%%%%%%%%%%%%%%%%%%%%%%%%%%%%%
%%
%% MARK: Covering Spaces
%%
%%%%%%%%%%%%%%%%%%%%%%%%%%%%%%%%%%%%%%%%%%%%%%%%%%%%%%%%%%%%%%%%%%%%%%%%%%%%%%%

%% \UniversalCover : Universal covering space (maximal simply connected cover)
%% Uses the stretching tilde from AMS.
%% Usage: \UniversalCover{\cT_{\alpha}} produces \widetilde{\cT_{\alpha}}
\newcommand{\UniversalCover}[1]{\widetilde{#1}}

%% \Cover : Any covering space (not necessarily universal)
%% Uses the hat from AMS. This is the general covering projection.
%% Usage: \Cover{G} produces \widehat{G}
\newcommand{\Cover}[1]{\widehat{#1}}

%%%%%%%%%%%%%%%%%%%%%%%%%%%%%%%%%%%%%%%%%%%%%%%%%%%%%%%%%%%%%%%%%%%%%%%%%%%%%%%
%%
%% MARK: Homotopy and Fundamental Group
%%
%%%%%%%%%%%%%%%%%%%%%%%%%%%%%%%%%%%%%%%%%%%%%%%%%%%%%%%%%%%%%%%%%%%%%%%%%%%%%%%

%% \FundamentalGroup : The fundamental group (first homotopy group)
%% Usage: \FundamentalGroup{X} produces π₁(X)
%% For a pointed space, use \FundamentalGroup{X,x₀}
\newcommand{\FundamentalGroup}[1]{\pi_1(#1)}

%% \piX : Alternative semantic name (shorter, still clear)
%% Usage: \piX{X} produces π₁(X)
\newcommand{\piX}[1]{\pi_1(#1)}

%%%%%%%%%%%%%%%%%%%%%%%%%%%%%%%%%%%%%%%%%%%%%%%%%%%%%%%%%%%%%%%%%%%%%%%%%%%%%%%
%%
%% MARK: Periods and Torus
%%
%%%%%%%%%%%%%%%%%%%%%%%%%%%%%%%%%%%%%%%%%%%%%%%%%%%%%%%%%%%%%%%%%%%%%%%%%%%%%%%

%% \Periods : Group of periods
\newcommand{\Periods}{\operatorname{P}}

%% \cTorus : Calligraphic T for distinctive infinite torus
\newcommand{\cTorus}{\cT}

%% \Torus : Standard torus R/Z
\newcommand{\Torus}{\mathrm{T}}

%% \Sph : n-sphere
\newcommand{\Sph}{\mathrm{S}}

%%%%%%%%%%%%%%%%%%%%%%%%%%%%%%%%%%%%%%%%%%%%%%%%%%%%%%%%%%%%%%%%%%%%%%%%%%%%%%%
%%
%% MARK: Linear Algebra and Vector Spaces
%%
%%%%%%%%%%%%%%%%%%%%%%%%%%%%%%%%%%%%%%%%%%%%%%%%%%%%%%%%%%%%%%%%%%%%%%%%%%%%%%%

%% \Lin : Space of linear maps between vector spaces
\newcommand{\Lin}{\operatorname{L}}

%% \DLin : Smooth linear maps between diffeological vector spaces
\newcommand{\DLin}{\Lin^\infty}

%% \Vector : Column vector with parentheses
%% Usage: \Vector{a \\ b \\ c}
\newcommand{\Vector}[1]{\begin{pmatrix}#1\end{pmatrix}}

%% \SqVect : Column vector with square brackets
\newcommand{\SqVect}[1]{\left[\begin{matrix}#1\end{matrix}\right]}

%% \Matrix : Set of matrices
\DeclareMathOperator{\Matrix}{Mat}

%% \ker : Kernel of a linear map
\renewcommand{\ker}{\mathrm{ker}}

%% \coker : Cokernel of a linear map
\newcommand{\coker}{\operatorname{coker}}

%% \J : The π/2 rotation matrix (complex structure)
\newcommand{\J}{\operatorname{J}}

%% \vecspan : Linear span of vectors
\newcommand{\vecspan}[1]{{\rm span}(#1)}

%%%%%%%%%%%%%%%%%%%%%%%%%%%%%%%%%%%%%%%%%%%%%%%%%%%%%%%%%%%%%%%%%%%%%%%%%%%%%%%
%%
%% MARK: Fonts (Semantic Variants)
%%
%% These provide semantic aliases for font variants.
%% The visual appearance is handled by Display-macros.tex
%%
%%%%%%%%%%%%%%%%%%%%%%%%%%%%%%%%%%%%%%%%%%%%%%%%%%%%%%%%%%%%%%%%%%%%%%%%%%%%%%%

%% Calligraphic letters (semantic: used for categories, sheaves, etc.)
\newcommand{\cA}{\mathcal{A}}
\newcommand{\cB}{\mathcal{B}}
\newcommand{\cC}{\mathcal{C}}
\newcommand{\cD}{\mathcal{D}}
\newcommand{\cE}{\mathcal{E}}
\newcommand{\cF}{\mathcal{F}}
\newcommand{\cG}{\mathcal{G}}
\newcommand{\cH}{\mathcal{H}}
\newcommand{\cI}{\mathcal{I}}
\newcommand{\cJ}{\mathcal{J}}
\newcommand{\cK}{\mathcal{K}}
\newcommand{\cL}{\mathcal{L}}
\newcommand{\cM}{\mathcal{M}}
\newcommand{\cN}{\mathcal{N}}
\newcommand{\cO}{\mathcal{O}}
\newcommand{\cP}{\mathcal{P}}
\newcommand{\cQ}{\mathcal{Q}}
\newcommand{\cR}{\mathcal{R}}
\newcommand{\cS}{\mathcal{S}}
\newcommand{\cT}{\mathcal{T}}
\newcommand{\cU}{\mathcal{U}}
\newcommand{\cV}{\mathcal{V}}
\newcommand{\cW}{\mathcal{W}}
\newcommand{\cX}{\mathcal{X}}
\newcommand{\cY}{\mathcal{Y}}
\newcommand{\cZ}{\mathcal{Z}}

%% Fraktur letters (semantic: used for Lie algebras, ideals, etc.)
\newcommand{\fA}{{\mathfrak A}}
\newcommand{\fB}{{\mathfrak B}}
\newcommand{\fC}{{\mathfrak C}}
\newcommand{\fD}{{\mathfrak D}}
\newcommand{\fF}{{\mathfrak F}}
\newcommand{\fG}{{\mathfrak G}}
\newcommand{\fK}{{\mathfrak K}}
\newcommand{\fI}{{\mathfrak I}}
\newcommand{\fL}{{\mathfrak L}}
\newcommand{\fP}{{\mathfrak P}}
\newcommand{\fR}{{\mathfrak R}}
\newcommand{\fS}{{\mathfrak S}}
\newcommand{\fU}{{\mathfrak U}}
\newcommand{\fX}{{\mathfrak X}}
\newcommand{\fY}{{\mathfrak Y}}
\newcommand{\fZ}{{\mathfrak Z}}
\newcommand{\fu}{{\mathfrak u}}

%% Bold letters (semantic: used for vectors, matrices, operators)
\renewcommand{\AA}{{\mathbf A}}  %% Works but loses the Angstrom symbol
\newcommand{\BB}{{\mathbf B}}
\newcommand{\DD}{{\mathbf D}}
\newcommand{\EE}{{\mathbf E}}
\newcommand{\FF}{{\mathbf F}}
\newcommand{\GG}{{\mathbf G}}
\newcommand{\HH}{{\mathbf H}}
\newcommand{\II}{{\mathbf I}}
\newcommand{\JJ}{{\mathbf J}}
\newcommand{\KK}{{\mathbf K}}
\newcommand{\LL}{{\mathbf L}}
\newcommand{\MM}{{\mathbf M}}
\newcommand{\OO}{{\mathbf O}}
\newcommand{\PP}{{\mathbf P}}
\renewcommand{\SS}{{\mathbf S}}
\newcommand{\TT}{{\mathbf T}}
\newcommand{\UU}{{\mathbf U}}
\newcommand{\VV}{{\mathbf V}}
\newcommand{\WW}{{\mathbf W}}
\newcommand{\XX}{{\mathbf X}}
\newcommand{\YY}{{\mathbf Y}}

\renewcommand{\aa}{{\bf a}}    %% Redefines "å" to bold a
\newcommand{\bb}{{\bf b}}
\newcommand{\ff}{{\bf f}}
\newcommand{\jj}{{\bf j}}
\newcommand{\rr}{{\bf r}}
\newcommand{\pp}{{\bf p}}
\renewcommand{\ss}{{\bf s}}
\newcommand{\uu}{{\bf u}}
\newcommand{\vv}{{\bf v}}
\newcommand{\xx}{{\bf x}}
\newcommand{\yy}{{\bf y}}
\newcommand{\ee}{{\bf e}}

%% Bold Greek letters
\newcommand{\balpha}{\boldsymbol \alpha}
\newcommand{\bbeta}{\boldsymbol \beta}
\newcommand{\bomega}{\boldsymbol \omega}
\newcommand{\bOmega}{\boldsymbol \Omega}
\newcommand{\bgamma}{\boldsymbol \gamma}
\newcommand{\bvarphi}{\boldsymbol \varphi}
\newcommand{\blambda}{\boldsymbol \lambda}
\newcommand{\bsigma}{\boldsymbol \sigma}
\newcommand{\bpi}{\boldsymbol \pi}
\newcommand{\btau}{\boldsymbol \tau}
\newcommand{\bPhi}{\boldsymbol{\Phi}}

%%%%%%%%%%%%%%%%%%%%%%%%%%%%%%%%%%%%%%%%%%%%%%%%%%%%%%%%%%%%%%%%%%%%%%%%%%%%%%%
%%
%% End of Diffeology-macros.tex
%%
%%%%%%%%%%%%%%%%%%%%%%%%%%%%%%%%%%%%%%%%%%%%%%%%%%%%%%%%%%%%%%%%%%%%%%%%%%%%%%%
